% Fichier engendré par tools/generate-tex.py à partir de 2025.
% Ne pas éditer à la main : tout le texte provient des commentaires du fichier Lean.
\documentclass[11pt,a4paper]{article}

% Compile aussi bien avec pdflatex qu'avec xelatex ou tectonic.
\usepackage{iftex}
\ifPDFTeX
  \usepackage[T1]{fontenc}
  \usepackage[utf8]{inputenc}
\else
  \usepackage{fontspec}
\fi
\usepackage[french]{babel}
\usepackage{amsmath,amssymb,amsthm}
\usepackage[margin=2.5cm]{geometry}
% Les figures : TikZ, pour que leurs étiquettes soient composées dans les
% polices du document. Voir la skill `illustrate-theorem`.
\usepackage{tikz}
\usetikzlibrary{calc}

\usepackage[hidelinks]{hyperref}

\theoremstyle{definition}
\newtheorem{definition}{Définition}
\newtheorem{exemple}{Exemple}
\theoremstyle{plain}
\newtheorem{theoreme}{Théorème}
\newtheorem{lemme}[theoreme]{Lemme}

% Un sujet pose des questions, non des théorèmes : voir exams/README.md.
\theoremstyle{definition}
\newtheorem*{reponse}{Réponse}
%% Ce qui remplace une démonstration quand il n'y en a pas : la déclaration Lean
%% porte un `sorry`, et ce qui suit dit pourquoi. Ce n'est pas une preuve.
\newtheorem*{enonceadmis}{Énoncé admis}

% Un exemple ne se démontre pas : on explique ce qu'il montre.
\newenvironment{explication}{\begin{proof}[Explication]}{\end{proof}}

% Un énoncé écrit mais non démontré : la preuve Lean est un `sorry`.
\newcommand{\admis}{\par\smallskip\noindent\textit{Admis : l'énoncé est écrit en Lean, sa démonstration reste à faire.}\par}

% Renvoi à la déclaration Lean, une ligne après la démonstration, aligné à droite.
\newcommand{\source}[2]{\par\smallskip\noindent\hfill{\footnotesize\href{#1}{\texttt{#2}}}\par\smallskip}

\title{Sujets 2025, décomposés}
\date{}

\begin{document}
\maketitle
% Les identifiants Lean en machine à écrire ne se coupent pas : on tolère des
% espacements plus lâches plutôt que des lignes qui débordent.
\sloppy

\section{Bac2025MetropoleJour1}

\noindent
Baccalauréat, épreuve de spécialité \textemdash{} voie générale, Métropole, 17 juin 2025 (jour 1).

Chaque question du sujet qui porte une proposition mathématique devient un théorème
Lean, dont la réponse est rédigée ci-dessous en français. Celles qui n'en portent pas \textemdash{} compléter un arbre pondéré, lire un nombre
dérivé sur un graphique, compléter un algorithme, interpréter un résultat \textemdash{} n'ont pas
d'énoncé, et l'index du dossier dit lesquelles et pourquoi.

Les probabilités sont traitées comme le sujet les traite : sur les nombres. La
construction de l'espace probabilisé n'apporterait rien à des calculs qui portent sur
des probabilités données, et le chapitre \og{} Probabilités et statistiques \fg{} du cours
fait ce travail-là pour ses propres énoncés.

\subsection{Exercice 1 \textemdash{} groupes sanguins, loi binomiale, Bienaymé-Tchebychev}

\begin{reponse}
Question 2 : \texttt{p(B \(\cap\) R) = p(B) \(\times\) p\_B(R) = 0,10 \(\times\) 0,84 = 0,084}.

C'est la lecture d'un chemin de l'arbre : les probabilités s'y multiplient.
\end{reponse}
\begin{proof}
On multiplie les probabilités portées par le chemin \og{} groupe $B$, puis rhésus
positif \fg{} :
\[ p(B \cap R) = 0{,}10 \times 0{,}84 = 0{,}084 . \]
\end{proof}
\source{https://github.com/Commutator-IO/learn-lean/blob/main/exams/2025/Bac2025MetropoleJour1.lean\#L24}{Bac2025MetropoleJour1.lean\#L24}

\begin{reponse}
Question 3 : \texttt{p\_O(R) = 0,83}.

Le groupe \texttt{O} est le complémentaire des trois autres, donc \texttt{p(O) = 1 − 0,45 − 0,10 −
0,03 = 0,42}. La formule des probabilités totales décompose \texttt{p(R)} sur les quatre
groupes : \texttt{p(O \(\cap\) R) = 0,839 7 − 0,45 \(\times\) 0,85 − 0,10 \(\times\) 0,84 − 0,03 \(\times\) 0,82 = 0,348 6}, et
il ne reste qu'à diviser par \texttt{p(O)}.
\end{reponse}

\begin{figure}[h]
\centering
% Construction : racine en (0, 0) ; les quatre groupes en x = 2,6, aux ordonnées
% 2,7 / 0,9 / −0,9 / −2,7 ; les deux rhésus en x = 5,4, décalés de ±0,55.
% La branche cherchée porte un point d'interrogation : c'est la question.
\begin{tikzpicture}[scale=0.72, font=\small]
  \coordinate (R) at (0, 0);
  \foreach \y/\g/\p in {2.7/A/0{,}45, 0.9/B/0{,}10, -0.9/AB/0{,}03, -2.7/O/0{,}42} {
    \draw (R) -- (2.6, \y);
    \node at (2.9, \y) {$\g$};
    \node[fill=white, inner sep=1pt] at (1.3, {\y / 2 + 0.22}) {$\p$};
  }
  \foreach \y/\p in {2.7/0{,}85, 0.9/0{,}84, -0.9/0{,}82} {
    \draw (3.25, \y) -- (5.4, {\y + 0.55});
    \draw (3.25, \y) -- (5.4, {\y - 0.55});
    \node[right] at (5.4, {\y + 0.55}) {$R$};
    \node[right] at (5.4, {\y - 0.55}) {$\bar{R}$};
    \node[fill=white, inner sep=1pt] at (4.4, {\y + 0.42}) {$\p$};
  }
  \draw (3.25, -2.7) -- (5.4, -2.15);
  \draw (3.25, -2.7) -- (5.4, -3.25);
  \node[right] at (5.4, -2.15) {$R$};
  \node[right] at (5.4, -3.25) {$\bar{R}$};
  \node[fill=white, inner sep=1pt] at (4.4, -2.28) {$?$};
\end{tikzpicture}

\caption{L'arbre pondéré des quatre groupes sanguins ; la branche cherchée porte le point d'interrogation.}
\end{figure}
\begin{proof}
Les quatre groupes forment une partition, donc leurs probabilités ont pour somme $1$ :
\[ p(O) = 1 - 0{,}45 - 0{,}10 - 0{,}03 = 0{,}42 . \]

La formule des probabilités totales décompose $p(R)$ sur cette partition,
\[ p(R) = p(A \cap R) + p(B \cap R) + p(AB \cap R) + p(O \cap R) , \]
d'où
\[ p(O \cap R) = 0{,}8397 - 0{,}45 \times 0{,}85 - 0{,}10 \times 0{,}84
   - 0{,}03 \times 0{,}82 = 0{,}8397 - 0{,}4911 = 0{,}3486 . \]

Enfin, par définition d'une probabilité conditionnelle,
\[ p_O(R) = \frac{p(O \cap R)}{p(O)} = \frac{0{,}3486}{0{,}42} = 0{,}83 . \]
\end{proof}
\source{https://github.com/Commutator-IO/learn-lean/blob/main/exams/2025/Bac2025MetropoleJour1.lean\#L32}{Bac2025MetropoleJour1.lean\#L32}

\begin{reponse}
Question 4 : la probabilité d'être donneur universel vaut \texttt{0,071 4}.

Un donneur universel est du groupe \texttt{O} et de rhésus négatif : \texttt{p(O \(\cap\) \(\overline{R}\)) = p(O) \(\times\)
(1 − p\_O(R)) = 0,42 \(\times\) 0,17}.
\end{reponse}
\begin{proof}
Un donneur universel est du groupe $O$ et de rhésus négatif. Les deux conditions se
composent le long d'un chemin de l'arbre :
\[ p(O \cap \bar{R}) = p(O) \times \bigl(1 - p_O(R)\bigr)
   = 0{,}42 \times 0{,}17 = 0{,}0714 . \]
\end{proof}
\source{https://github.com/Commutator-IO/learn-lean/blob/main/exams/2025/Bac2025MetropoleJour1.lean\#L41}{Bac2025MetropoleJour1.lean\#L41}

\begin{reponse}
Question 5c : \texttt{E(X) = 7,14} et \texttt{V(X) = 6,63} à \texttt{10\(^{-2}\)} près.

Pour une loi binomiale de paramètres \texttt{n} et \texttt{p}, l'espérance vaut \texttt{np} et la variance
\texttt{np(1 − p)} : ce sont les deux formules du cours, et le sujet ne demande que de les
appliquer. La variance exacte est \texttt{6,629 604}, dont l'arrondi au centième est \texttt{6,63}.
\end{reponse}
\begin{proof}
Pour une loi binomiale de paramètres $n$ et $p$, l'espérance vaut $np$ et la variance
$np(1-p)$. Ici $n = 100$ et $p = 0{,}0714$, donc
\[ E(X) = 100 \times 0{,}0714 = 7{,}14 \]
et
\[ V(X) = 100 \times 0{,}0714 \times 0{,}9286 = 6{,}629604 . \]

L'arrondi au centième se lit sur l'écart à $6{,}63$ :
\[ |6{,}629604 - 6{,}63| = 0{,}000396 < 0{,}005 , \]
ce qui est exactement la définition de l'arrondi au centième.
\end{proof}
\source{https://github.com/Commutator-IO/learn-lean/blob/main/exams/2025/Bac2025MetropoleJour1.lean\#L48}{Bac2025MetropoleJour1.lean\#L48}

\begin{reponse}
Question 6b : \texttt{E(M\_N) = 7,14}, quel que soit le nombre \texttt{N} de villes.

L'espérance est linéaire : celle d'une moyenne de variables de même espérance est
cette espérance. Grouper les villes ne change donc rien à ce qu'on attend en
moyenne.
\end{reponse}
\begin{proof}
La moyenne des $N$ variables a pour espérance $\dfrac{1}{N}$ fois la somme de leurs
espérances, toutes égales à $7{,}14$ :
\[ \frac{N \times 7{,}14}{N} = 7{,}14 . \]
La simplification demande $N \ne 0$, ce que l'hypothèse fournit \textemdash{} et qui va de
soi, une moyenne sur zéro ville n'ayant pas de sens.
\end{proof}
\source{https://github.com/Commutator-IO/learn-lean/blob/main/exams/2025/Bac2025MetropoleJour1.lean\#L59}{Bac2025MetropoleJour1.lean\#L59}

\begin{reponse}
Question 6c : \texttt{V(M\_N) = 6,63/N}.

Les variables étant indépendantes, leurs variances s'ajoutent ; la division par \texttt{N}
étant une multiplication par \texttt{1/N}, elle divise la variance par \texttt{N\(^{2}\)}. C'est là que
l'hypothèse d'indépendance sert, et nulle part ailleurs \textemdash{} et c'est ce quotient par
\texttt{N} qui fait que la moyenne se resserre quand on multiplie les villes.
\end{reponse}
\begin{proof}
Les $N$ variables étant indépendantes, leurs variances s'ajoutent : la somme a pour
variance $N \times 6{,}63$. Diviser par $N$, c'est multiplier par $\dfrac{1}{N}$, ce qui
multiplie la variance par $\dfrac{1}{N^2}$. D'où
\[ V(M_N) = \frac{N \times 6{,}63}{N^2} = \frac{6{,}63}{N} , \]
la simplification étant licite puisque $N \ne 0$.

C'est ce quotient par $N$ qui fait que la moyenne se resserre quand on multiplie les
villes ; l'espérance, elle, ne bouge pas.
\end{proof}
\source{https://github.com/Commutator-IO/learn-lean/blob/main/exams/2025/Bac2025MetropoleJour1.lean\#L68}{Bac2025MetropoleJour1.lean\#L68}

\begin{reponse}
Question 6d : la plus petite valeur de \texttt{N} convenable est \texttt{6 766}.

L'inégalité de Bienaymé-Tchebychev majore \texttt{P(|M\_N − 7,14| \(\ge\) 0,14)} par
\texttt{V(M\_N)/0,14\(^{2}\) = 6,63/(0,019 6 N)}. Pour que l'évènement contraire \textemdash{} c'est-à-dire
\texttt{7 < M\_N < 7,28} \textemdash{} ait une probabilité au moins \texttt{0,95}, il suffit que cette majoration
descende à \texttt{0,05}, soit \texttt{N \(\ge\) 6 765,3\(\dots\)}.

L'inégalité obtenue sur les rationnels est stricte à \texttt{6 765} et vraie à \texttt{6 766} : le
seuil n'est pas entier, et c'est le passage à l'entier supérieur qui donne la
réponse.
\end{reponse}
\begin{proof}
Comme $N$ est un entier strictement positif, $N > 0$ et $N \times 0{,}14^2 > 0$ : on peut
multiplier les deux membres par cette quantité sans changer le sens de l'inégalité. La
condition s'écrit donc
\[ 6{,}63 \le 0{,}05 \times 0{,}0196 \times N = 0{,}00098\,N . \]

\emph{Sens direct.} Si cette inégalité a lieu, alors $N \ge \dfrac{6{,}63}{0{,}00098}
= 6\,765{,}30\ldots$, donc $N > 6\,765$ ; un entier strictement supérieur à $6\,765$ vaut
au moins $6\,766$.

\emph{Réciproque.} Si $N \ge 6\,766$, alors $0{,}00098 \times N \ge 0{,}00098 \times
6\,766 = 6{,}63068 \ge 6{,}63$.

Le seuil exact n'étant pas entier, c'est le passage à l'entier supérieur qui donne la
réponse $6\,766$.
\end{proof}
\source{https://github.com/Commutator-IO/learn-lean/blob/main/exams/2025/Bac2025MetropoleJour1.lean\#L81}{Bac2025MetropoleJour1.lean\#L81}

\subsection{Exercice 2 \textemdash{} une fonction avec un logarithme}

\begin{definition}
La fonction de l'exercice : \texttt{f(x) = x[2(ln x)\(^{2}\) − 3 ln x + 2]}, sur \texttt{]0 ; +\(\infty\)[}.
\end{definition}
\source{https://github.com/Commutator-IO/learn-lean/blob/main/exams/2025/Bac2025MetropoleJour1.lean\#L97}{Bac2025MetropoleJour1.lean\#L97}

\begin{reponse}
Partie B, question 1 : l'équation \texttt{2X\(^{2}\) − 3X + 2 = 0} n'a pas de solution réelle,
et la courbe de \texttt{f} ne coupe donc pas l'axe des abscisses.

Le discriminant vaut \texttt{9 − 16 = −7 < 0}. Ce qui se démontre le plus directement est
la mise sous forme canonique : \texttt{2X\(^{2}\) − 3X + 2 = 2(X − 3/4)\(^{2}\) + 7/8 \(\ge\) 7/8 > 0}. Le
crochet ne s'annulant pas et \texttt{x} étant strictement positif, \texttt{f(x)} ne s'annule pas
non plus.
\end{reponse}
\begin{proof}
On met le trinôme sous forme canonique :
\[ 2\left(X - \frac{3}{4}\right)^2 + \frac{7}{8}
   = 2X^2 - 3X + \frac{9}{8} + \frac{7}{8} = 2X^2 - 3X + 2 . \]
Un carré est positif ou nul, donc $2\bigl(X - \tfrac{3}{4}\bigr)^2 \ge 0$ et
\[ 2X^2 - 3X + 2 \ge \frac{7}{8} . \]

C'est la traduction du discriminant $9 - 16 = -7 < 0$, sous une forme qui donne en prime
une minoration explicite.
\end{proof}
\source{https://github.com/Commutator-IO/learn-lean/blob/main/exams/2025/Bac2025MetropoleJour1.lean\#L106}{Bac2025MetropoleJour1.lean\#L106}

\begin{reponse}
Le corollaire : \texttt{f} ne s'annule sur \texttt{]0 ; +\(\infty\)[}.
\end{reponse}
\begin{proof}
Appliquons la minoration précédente à $X = \ln x$ : le crochet
$2(\ln x)^2 - 3\ln x + 2$ vaut au moins $\dfrac{7}{8}$, donc il est strictement positif.

Par définition, $f(x)$ est le produit de $x$, strictement positif par hypothèse, et de ce
crochet. Un produit de deux facteurs strictement positifs est strictement positif, donc
$f(x) > 0$ et en particulier $f(x) \ne 0$.
\end{proof}
\source{https://github.com/Commutator-IO/learn-lean/blob/main/exams/2025/Bac2025MetropoleJour1.lean\#L110}{Bac2025MetropoleJour1.lean\#L110}

\begin{reponse}
Partie B, question 2 : \texttt{f} tend vers \texttt{+\(\infty\)} en \texttt{+\(\infty\)}.

Le sujet passe par la forme \texttt{x[2(ln x)\(^{2}\) − 3 ln x + 2]} et le fait que le crochet tend
vers \texttt{+\(\infty\)}. La minoration de la question précédente suffit et dit davantage : le
crochet reste au-dessus de \texttt{7/8}, donc \texttt{f(x) \(\ge\) 7x/8}, et une fonction minorée par une
fonction qui tend vers \texttt{+\(\infty\)} tend vers \texttt{+\(\infty\)}.
\end{reponse}
\begin{proof}
Pour tout $x > 0$, la minoration du crochet donne, en multipliant par $x \ge 0$,
\[ f(x) = \bigl(2(\ln x)^2 - 3\ln x + 2\bigr) \, x \ge \frac{7}{8}\,x . \]

Or $x \mapsto \dfrac{7}{8}x$ tend vers $+\infty$ en $+\infty$, comme produit d'une
constante strictement positive par l'identité. Une fonction minorée à partir d'un certain
rang par une fonction qui tend vers $+\infty$ tend elle-même vers $+\infty$ : c'est le
théorème de comparaison, et il conclut.
\end{proof}
\source{https://github.com/Commutator-IO/learn-lean/blob/main/exams/2025/Bac2025MetropoleJour1.lean\#L123}{Bac2025MetropoleJour1.lean\#L123}

\begin{reponse}
Partie B, question 3a : \texttt{f''(x) = (4 ln x + 1)/x} sur \texttt{]0 ; +\(\infty\)[}.

Le sujet admet \texttt{f'(x) = 2(ln x)\(^{2}\) + ln x − 1} ; la dérivée seconde s'en déduit en
dérivant terme à terme, \texttt{(ln x)\(^{2}\)} donnant \texttt{2 ln x / x}.
\end{reponse}
\begin{proof}
Sur $]0 ; +\infty[$, la dérivée du logarithme est $\dfrac{1}{x}$. On dérive terme à
terme la fonction admise $f'(t) = 2(\ln t)^2 + \ln t - 1$ :

\begin{itemize}
\item $(\ln t)^2$ a pour dérivée $2 \ln x \times \dfrac{1}{x}$, par dérivation d'une
      composée ;
\item $\ln t$ a pour dérivée $\dfrac{1}{x}$ ;
\item la constante $-1$ a une dérivée nulle.
\end{itemize}

En sommant,
\[ f''(x) = 2 \times \frac{2 \ln x}{x} + \frac{1}{x}
   = \frac{4 \ln x + 1}{x} . \]
\end{proof}
\source{https://github.com/Commutator-IO/learn-lean/blob/main/exams/2025/Bac2025MetropoleJour1.lean\#L140}{Bac2025MetropoleJour1.lean\#L140}

\begin{reponse}
Partie B, question 3b : la dérivée seconde est positive si et seulement si
\texttt{x \(\ge\) e\textasciicircum{}(−1/4)}.

C'est le signe de \texttt{4 ln x + 1}, le dénominateur \texttt{x} étant positif. La fonction est
donc concave avant \texttt{e\textasciicircum{}(−1/4)}, convexe après, et ce point est son point d'inflexion :
l'abscisse exacte est \texttt{e\textasciicircum{}(−1/4)}, non une valeur décimale.

Ce qui est démontré est le signe ; la convexité s'en déduit par le critère du cours,
qui n'est pas repris ici.
\end{reponse}
\begin{proof}
Le dénominateur $x$ étant strictement positif, l'inégalité $0 \le \dfrac{4 \ln x + 1}{x}$
équivaut à $0 \le 4 \ln x + 1$, c'est-à-dire à
\[ \ln x \ge -\frac{1}{4} . \]

Le logarithme et l'exponentielle sont deux bijections réciproques croissantes : pour
$x > 0$, l'inégalité $-\tfrac{1}{4} \le \ln x$ équivaut à $\mathrm{e}^{-1/4} \le x$.

C'est le signe de $f''$ qui est établi ici ; la concavité avant
$\mathrm{e}^{-1/4}$ et la convexité après s'en déduisent par le critère du cours, qui
n'est pas repris.
\end{proof}
\source{https://github.com/Commutator-IO/learn-lean/blob/main/exams/2025/Bac2025MetropoleJour1.lean\#L158}{Bac2025MetropoleJour1.lean\#L158}

\begin{reponse}
Partie C, question 1 : la tangente à la courbe au point \texttt{B(e ; e)} a pour
équation réduite \texttt{y = 2x − e}.

En \texttt{x = e}, le logarithme vaut \texttt{1} : \texttt{f(e) = e(2 − 3 + 2) = e} et
\texttt{f'(e) = 2 + 1 − 1 = 2}. L'équation de la tangente est
\texttt{y = f'(e)(x − e) + f(e) = 2x − e}.
\end{reponse}
\begin{proof}
En $x = \mathrm{e}$, le logarithme vaut $1$ puisque logarithme et exponentielle sont
réciproques l'un de l'autre. Donc
\[ f(\mathrm{e}) = \mathrm{e}\,(2 - 3 + 2) = \mathrm{e} \]
et
\[ f'(\mathrm{e}) = 2 \times 1^2 + 1 - 1 = 2 . \]

L'équation réduite de la tangente au point d'abscisse $\mathrm{e}$ est alors
\[ y = f'(\mathrm{e})(x - \mathrm{e}) + f(\mathrm{e}) = 2x - 2\mathrm{e} + \mathrm{e}
   = 2x - \mathrm{e} . \]
\end{proof}
\source{https://github.com/Commutator-IO/learn-lean/blob/main/exams/2025/Bac2025MetropoleJour1.lean\#L169}{Bac2025MetropoleJour1.lean\#L169}

\begin{reponse}
Partie C, question 3 : l'aire du domaine hachuré vaut \texttt{(3e\(^{2}\) − 4e − 5)/4} unités
d'aire.

Le domaine est compris entre la courbe et sa tangente, entre \texttt{x = 1} et \texttt{x = e} :
son aire est \texttt{\(\int\)\(_{1}\)\textasciicircum{}e (f(x) − (2x − e)) dx}, et l'intégrande se simplifie en
\texttt{2x(ln x)\(^{2}\) − 3x ln x + e}. Les deux intégrales sont celles de la question 2 et de
l'énoncé \textemdash{} \texttt{(e\(^{2}\) + 1)/4} et \texttt{(e\(^{2}\) − 1)/4} \textemdash{} et sont ici prises en hypothèse, comme le
sujet les admet. Reste \texttt{2(e\(^{2}\)−1)/4 − 3(e\(^{2}\)+1)/4 + e(e−1)}.
\end{reponse}

\begin{figure}[h]
\centering
% Construction : f(x) = x[2(ln x)² − 3 ln x + 2], et sa tangente y = 2x − e au
% point B(e ; e). Le domaine hachuré est compris entre la courbe et la tangente,
% de x = 1 à x = e ; son aire vaut (3e² − 4e − 5)/4.
\begin{tikzpicture}[xscale=1.7, yscale=0.85]
  \draw[->] (0, 0) -- (3.15, 0) node[below] {$x$};
  \draw[->] (0, -1) -- (0, 3.15);
  \draw[thin] (1.1322, -0.4539) -- (1.1322, 1.8776);
  \draw[thin] (1.2644, -0.1896) -- (1.2644, 1.7781);
  \draw[thin] (1.3965, 0.0748) -- (1.3965, 1.7053);
  \draw[thin] (1.5287, 0.3391) -- (1.5287, 1.6617);
  \draw[thin] (1.6609, 0.6035) -- (1.6609, 1.6489);
  \draw[thin] (1.7931, 0.8678) -- (1.7931, 1.6678);
  \draw[thin] (1.9252, 1.1322) -- (1.9252, 1.7193);
  \draw[thin] (2.0574, 1.3965) -- (2.0574, 1.8036);
  \draw[thin] (2.1896, 1.6609) -- (2.1896, 1.9209);
  \draw[thin] (2.3218, 1.9252) -- (2.3218, 2.0711);
  \draw[thin] (2.4539, 2.1896) -- (2.4539, 2.2542);
  \draw[thin] (2.5861, 2.4539) -- (2.5861, 2.4700);
  \draw[thick] plot coordinates {(1.0000, 2.0000) (1.0430, 1.9580) (1.0859, 1.9181) (1.1289, 1.8804) (1.1718, 1.8451) (1.2148, 1.8125) (1.2577, 1.7825) (1.3007, 1.7553) (1.3437, 1.7311) (1.3866, 1.7098) (1.4296, 1.6916) (1.4725, 1.6766) (1.5155, 1.6647) (1.5584, 1.6561) (1.6014, 1.6508) (1.6444, 1.6487) (1.6873, 1.6501) (1.7303, 1.6548) (1.7732, 1.6629) (1.8162, 1.6745) (1.8591, 1.6895) (1.9021, 1.7079) (1.9451, 1.7298) (1.9880, 1.7552) (2.0310, 1.7841) (2.0739, 1.8164) (2.1169, 1.8523) (2.1598, 1.8916) (2.2028, 1.9344) (2.2458, 1.9807) (2.2887, 2.0305) (2.3317, 2.0837) (2.3746, 2.1404) (2.4176, 2.2006) (2.4605, 2.2643) (2.5035, 2.3314) (2.5465, 2.4019) (2.5894, 2.4759) (2.6324, 2.5533) (2.6753, 2.6341) (2.7183, 2.7183)};
  \draw (1, {2 - 2.71828}) -- (3, {6 - 2.71828});
  \fill (2.71828, 2.71828) circle (0.9pt);
  \node[above left] at (2.71828, 2.71828) {$B$};
  \node[below] at (1, -0.06) {$1$};
  \node[below] at (2.71828, -0.06) {$\mathrm{e}$};
  \node[right] at (2.6, 0.75) {$\mathcal{C}_f$};
  \node[right] at (3, 3.28) {$y = 2x - \mathrm{e}$};
\end{tikzpicture}

\caption{Le domaine compris entre la courbe de \(f\) et sa tangente en \(B\), de \(x = 1\) à \(x = \mathrm{e}\).}
\end{figure}
\begin{proof}
On remplace $I$ et $J$ par les valeurs admises et l'on réduit au même dénominateur :
\[
\begin{aligned}
2J - 3I + \mathrm{e}(\mathrm{e} - 1)
 &= \frac{2(\mathrm{e}^2 - 1)}{4} - \frac{3(\mathrm{e}^2 + 1)}{4}
    + \mathrm{e}^2 - \mathrm{e} \\
 &= \frac{2\mathrm{e}^2 - 2 - 3\mathrm{e}^2 - 3}{4} + \mathrm{e}^2 - \mathrm{e} \\
 &= \frac{-\mathrm{e}^2 - 5 + 4\mathrm{e}^2 - 4\mathrm{e}}{4}
  = \frac{3\mathrm{e}^2 - 4\mathrm{e} - 5}{4} .
\end{aligned}
\]

Les deux intégrales $I$ et $J$ sont ici prises en hypothèse, comme le sujet les admet :
ce qui est démontré est la combinaison, non leur calcul.
\end{proof}
\source{https://github.com/Commutator-IO/learn-lean/blob/main/exams/2025/Bac2025MetropoleJour1.lean\#L182}{Bac2025MetropoleJour1.lean\#L182}

\subsection{Exercice 3 \textemdash{} vrai ou faux dans l'espace}

\begin{reponse}
Affirmation 1, \textbf{vraie}. Le vecteur directeur proposé est \texttt{(−2 ; −1 ; 3)}, et
\texttt{AB = (4 ; 2 ; −6) = −2 \(\times\) (−2 ; −1 ; 3)} : il est colinéaire à \texttt{AB}. Le point obtenu
pour \texttt{t = 0} est \texttt{B(3 ; 2 ; −1)}, qui appartient à la droite. La représentation
paramétrique est donc bien celle de \texttt{(AB)}.
\end{reponse}
\begin{proof}
On calcule les coordonnées du vecteur $\vec{AB}$, différence des coordonnées de $B$ et de
$A$ :
\[ \vec{AB} = (3 - (-1) ;\ 2 - 0 ;\ -1 - 5) = (4 ;\ 2 ;\ -6) , \]
et
\[ -2 \times (-2 ;\ -1 ;\ 3) = (4 ;\ 2 ;\ -6) . \]

Les deux triplets coïncident : le vecteur directeur proposé est colinéaire à $\vec{AB}$.
Comme la valeur $t = 0$ du paramètre donne le point $B$, la représentation paramétrique
est bien celle de la droite $(AB)$.
\end{proof}
\source{https://github.com/Commutator-IO/learn-lean/blob/main/exams/2025/Bac2025MetropoleJour1.lean\#L194}{Bac2025MetropoleJour1.lean\#L194}

\begin{reponse}
Affirmation 2, \textbf{fausse}. Le vecteur \texttt{n = (5 ; −2 ; 1)} est bien orthogonal à
\texttt{OA} \textemdash{} le produit scalaire vaut \texttt{−5 + 0 + 5 = 0} \textemdash{} mais pas à \texttt{OB}, où il vaut
\texttt{15 − 4 − 1 = 10}. Un vecteur normal à un plan doit être orthogonal à deux vecteurs
non colinéaires du plan : une seule orthogonalité ne suffit pas, et c'est le piège de
l'affirmation.
\end{reponse}
\begin{proof}
Le produit scalaire est, dans un repère orthonormé, la somme des produits des coordonnées.
Pour $\vec{n}(5 ; -2 ; 1)$ et $\vec{OA}(-1 ; 0 ; 5)$ :
\[ 5 \times (-1) + (-2) \times 0 + 1 \times 5 = -5 + 0 + 5 = 0 . \]
Pour $\vec{n}$ et $\vec{OB}(3 ; 2 ; -1)$ :
\[ 5 \times 3 + (-2) \times 2 + 1 \times (-1) = 15 - 4 - 1 = 10 \ne 0 . \]

Un vecteur normal à un plan doit être orthogonal à deux vecteurs non colinéaires de ce
plan : la première orthogonalité ne suffit pas, et la seconde est en défaut. L'affirmation
est fausse.
\end{proof}
\source{https://github.com/Commutator-IO/learn-lean/blob/main/exams/2025/Bac2025MetropoleJour1.lean\#L203}{Bac2025MetropoleJour1.lean\#L203}

\begin{reponse}
Affirmation 3, \textbf{fausse} : les droites \texttt{d} et \texttt{d'} sont sécantes, donc
coplanaires.

Le système donné par les deux premières coordonnées a pour solution \texttt{s = 5/2} et
\texttt{k = −4}, et la troisième coordonnée est alors satisfaite aussi : les deux droites
passent par \texttt{(−9 ; 12 ; −14)}. Démontrer qu'une affirmation d'existence est fausse
c'est exhiber le point, et c'est plus court que de raisonner sur les vecteurs.
\end{reponse}
\begin{proof}
Démontrer qu'une affirmation de non-existence est fausse, c'est exhiber un contre-exemple.
Prenons $k = -4$ et $s = \dfrac{5}{2}$, et vérifions les trois coordonnées :
\[ 15 + (-4) = 11 = 1 + 4 \times \frac{5}{2} , \qquad
   8 - (-4) = 12 = 2 + 4 \times \frac{5}{2} , \]
\[ -6 + 2 \times (-4) = -14 = 1 - 6 \times \frac{5}{2} . \]

Les deux droites passent donc toutes deux par le point $(-9 ; 12 ; -14)$ : elles sont
sécantes, donc coplanaires, et l'affirmation est fausse.

Les deux premières équations suffisaient à déterminer le couple $(k, s)$ ; c'est la
troisième, satisfaite \og{} par surcroît \fg{}, qui fait la conclusion.
\end{proof}
\source{https://github.com/Commutator-IO/learn-lean/blob/main/exams/2025/Bac2025MetropoleJour1.lean\#L214}{Bac2025MetropoleJour1.lean\#L214}

\begin{reponse}
Affirmation 4, \textbf{vraie} : la distance de \texttt{C(2 ; −1 ; 2)} au plan
\texttt{x − y + z + 1 = 0} vaut \texttt{2\(\sqrt{3}\)}.

La formule donne \texttt{|2 + 1 + 2 + 1|/\(\sqrt{1 + 1 + 1}\) = 6/\(\sqrt{3}\)}, et \texttt{6/\(\sqrt{3}\) = 2\(\sqrt{3}\)} \textemdash{} c'est la
même valeur écrite sans radical au dénominateur.
\end{reponse}
\begin{proof}
La distance d'un point à un plan d'équation $ax + by + cz + d = 0$ est
$\dfrac{|ax_0 + by_0 + cz_0 + d|}{\sqrt{a^2 + b^2 + c^2}}$. Ici le numérateur vaut
\[ |2 - (-1) + 2 + 1| = |6| = 6 , \]
et le dénominateur $\sqrt{1 + 1 + 1} = \sqrt{3}$.

Reste à voir que $\dfrac{6}{\sqrt{3}} = 2\sqrt{3}$. Comme $\sqrt{3} > 0$, il revient au
même de vérifier, après multiplication par $\sqrt{3}$,
\[ 6 = 2\sqrt{3} \times \sqrt{3} = 2 \times 3 , \]
où l'on a utilisé $\sqrt{3} \times \sqrt{3} = 3$ \textemdash{} la racine carrée d'un réel
positif est le réel positif dont le carré vaut ce réel.
\end{proof}
\source{https://github.com/Commutator-IO/learn-lean/blob/main/exams/2025/Bac2025MetropoleJour1.lean\#L223}{Bac2025MetropoleJour1.lean\#L223}

\subsection{Exercice 4 \textemdash{} la posidonie}

\begin{definition}
La suite du modèle discret : \texttt{u\(_{0}\) = 1} et \texttt{u\_\{n+1\} = −0,02u\_n\(^{2}\) + 1,3u\_n}.
\end{definition}
\source{https://github.com/Commutator-IO/learn-lean/blob/main/exams/2025/Bac2025MetropoleJour1.lean\#L234}{Bac2025MetropoleJour1.lean\#L234}

\begin{reponse}
La relation de récurrence, sous la forme où les preuves l'utilisent.
\end{reponse}
\begin{proof}
Les deux membres sont le même énoncé : c'est la seconde clause de la définition récursive
de $u$, celle qui donne $u_{n+1}$ à partir de $u_n$.
\end{proof}
\source{https://github.com/Commutator-IO/learn-lean/blob/main/exams/2025/Bac2025MetropoleJour1.lean\#L239}{Bac2025MetropoleJour1.lean\#L239}

\begin{reponse}
Partie A, question 1 : au premier juillet 2025, la posidonie recouvre
\texttt{1,28 hectare}.
\end{reponse}
\begin{proof}
On applique la relation de récurrence à $n = 0$, avec $u_0 = 1$ :
\[ u_1 = -0{,}02 \times 1^2 + 1{,}3 \times 1 = -0{,}02 + 1{,}3 = 1{,}28 . \]
\end{proof}
\source{https://github.com/Commutator-IO/learn-lean/blob/main/exams/2025/Bac2025MetropoleJour1.lean\#L243}{Bac2025MetropoleJour1.lean\#L243}

\begin{reponse}
Partie A, question 2a : pour tout \texttt{n}, \texttt{1 \(\le\) u\_n \(\le\) u\_\{n+1\} \(\le\) 20}.

La récurrence porte sur les trois inégalités à la fois, et c'est nécessaire : la
croissance de \texttt{n} à \texttt{n+1} se déduit de celle de \texttt{n−1} à \texttt{n} en appliquant \texttt{h}, non
d'un raisonnement sur le point fixe \textemdash{} la fonction \texttt{h(x) = −0,02x\(^{2}\) + 1,3x} est
croissante sur \texttt{[0 ; 20]}, ce que le sujet admet et que les inégalités polynomiales
retrouvent directement. Sans le majorant \texttt{20}, cette monotonie tomberait : \texttt{h}
décroît au-delà de \texttt{32,5}.
\end{reponse}

\begin{figure}[h]
\centering
% Construction : h(x) = −0,02x² + 1,3x, la première bissectrice y = x, et
% l'escalier parti de u₀ = 1. Le point fixe est en (15 ; 15) : la suite croît
% en restant sous 20, et s'en approche sans l'atteindre.
\begin{tikzpicture}[scale=0.29]
  \draw[->] (0, 0) -- (18.5, 0) node[below] {$x$};
  \draw[->] (0, 0) -- (0, 18.5);
  \draw[dashed] (0, 0) -- (17.5, 17.5) node[above left] {$y = x$};
  \draw[thick] plot coordinates {(0.000, 0.000) (0.500, 0.645) (1.000, 1.280) (1.500, 1.905) (2.000, 2.520) (2.500, 3.125) (3.000, 3.720) (3.500, 4.305) (4.000, 4.880) (4.500, 5.445) (5.000, 6.000) (5.500, 6.545) (6.000, 7.080) (6.500, 7.605) (7.000, 8.120) (7.500, 8.625) (8.000, 9.120) (8.500, 9.605) (9.000, 10.080) (9.500, 10.545) (10.000, 11.000) (10.500, 11.445) (11.000, 11.880) (11.500, 12.305) (12.000, 12.720) (12.500, 13.125) (13.000, 13.520) (13.500, 13.905) (14.000, 14.280) (14.500, 14.645) (15.000, 15.000) (15.500, 15.345) (16.000, 15.680) (16.500, 16.005) (17.000, 16.320)};
  \draw (1.000, 1.000) -- (1.000, 1.280) -- (1.280, 1.280) -- (1.280, 1.631) -- (1.631, 1.631) -- (1.631, 2.067) -- (2.067, 2.067) -- (2.067, 2.602) -- (2.602, 2.602) -- (2.602, 3.247) -- (3.247, 3.247) -- (3.247, 4.011) -- (4.011, 4.011) -- (4.011, 4.892) -- (4.892, 4.892) -- (4.892, 5.881) -- (5.881, 5.881) -- (5.881, 6.954) -- (6.954, 6.954) -- (6.954, 8.073) -- (8.073, 8.073) -- (8.073, 9.191) -- (9.191, 9.191) -- (9.191, 10.259) -- (10.259, 10.259) -- (10.259, 11.232) -- (11.232, 11.232) -- (11.232, 12.078) -- (12.078, 12.078);
  \fill (15, 15) circle (5pt);
  \node[right] at (15.4, 14.4) {$15$};
  \node[below] at (1, -0.3) {$u_0$};
  \node[right] at (4.2, 13.6) {$y = h(x)$};
\end{tikzpicture}

\caption{L'escalier de la suite : partie de \(u_0 = 1\), elle croît en restant sous \(20\), vers le point fixe \(15\).}
\end{figure}
\begin{proof}
Posons $h(x) = -0{,}02x^2 + 1{,}3x$, de sorte que $u_{n+1} = h(u_n)$. Deux propriétés de
$h$ servent, et il faut les établir avant la récurrence.

\emph{$h$ est croissante sur $[0 ; 20]$.} Pour $0 \le a \le b \le 20$,
\[ h(b) - h(a) = (b - a)\bigl(1{,}3 - 0{,}02\,(a + b)\bigr) . \]
Le premier facteur est positif ou nul ; le second aussi, puisque $a + b \le 40$ donne
$0{,}02\,(a+b) \le 0{,}8 \le 1{,}3$. Donc $h(a) \le h(b)$.

\emph{$h$ ne dépasse pas $20$ sur $[1 ; 20]$.} Pour $1 \le b \le 20$,
\[ 20 - h(b) = 0{,}02\,b^2 - 1{,}3\,b + 20 = 0{,}02\,(25 - b)(40 - b) , \]
produit de deux facteurs positifs puisque $b \le 20 < 25 < 40$. Donc $h(b) \le 20$.

\emph{Récurrence.} Au rang $0$ : $u_0 = 1$, donc $1 \le u_0$ ; et $u_1 = 1{,}28$, donc
$u_0 \le u_1$ et $u_1 \le 20$.

Supposons $1 \le u_n \le u_{n+1} \le 20$. Alors $1 \le u_{n+1}$ par transitivité. Ensuite,
la croissance de $h$ appliquée à $0 \le u_n \le u_{n+1} \le 20$ donne
\[ u_{n+1} = h(u_n) \le h(u_{n+1}) = u_{n+2} , \]
et la majoration appliquée à $1 \le u_{n+1} \le 20$ donne $u_{n+2} = h(u_{n+1}) \le 20$.
Les trois inégalités sont vraies au rang $n+1$.

Il fallait porter la récurrence sur les trois inégalités à la fois : la croissance d'un
rang au suivant se déduit de celle du rang précédent, et le majorant $20$ est ce qui
maintient la suite dans la zone où $h$ est croissante \textemdash{} $h$ décroît au-delà de
$32{,}5$.
\end{proof}
\source{https://github.com/Commutator-IO/learn-lean/blob/main/exams/2025/Bac2025MetropoleJour1.lean\#L255}{Bac2025MetropoleJour1.lean\#L255}

\begin{reponse}
Partie A, question 2c : la limite de la suite vaut \texttt{15 hectares}.

Si la suite converge vers \texttt{L}, la limite vérifie \texttt{L = −0,02L\(^{2}\) + 1,3L}, c'est-à-dire
\texttt{0,02L\(^{2}\) = 0,3L} : les points fixes sont \texttt{0} et \texttt{15}. La suite étant minorée par \texttt{1},
sa limite l'est aussi, et \texttt{L = 15}. La zone étudiée fait \texttt{20 hectares} : le modèle
prévoit donc que la posidonie n'en recouvrira jamais toute la surface.
\end{reponse}
\begin{proof}
Si la suite converge vers $L$, alors $L$ est un point fixe de la relation de récurrence :
$L = -0{,}02L^2 + 1{,}3L$. En regroupant, $0{,}02L^2 = 0{,}3L$, c'est-à-dire, après
division par $0{,}02$,
\[ L(L - 15) = 0 . \]

Un produit est nul quand l'un de ses facteurs l'est. Le cas $L = 0$ est exclu par
l'hypothèse $L \ge 1$ ; reste $L = 15$.

La zone étudiée faisant $20$ hectares, le modèle prévoit que la posidonie n'en recouvrira
jamais toute la surface.
\end{proof}
\source{https://github.com/Commutator-IO/learn-lean/blob/main/exams/2025/Bac2025MetropoleJour1.lean\#L283}{Bac2025MetropoleJour1.lean\#L283}

\begin{reponse}
Partie B, question 1 : \texttt{g = 1/f} est solution de \texttt{y' = −0,3y + 0,02}.

C'est le changement de fonction qui rend l'équation linéaire : la dérivée de \texttt{1/f}
vaut \texttt{−f'/f\(^{2}\)}, et en remplaçant \texttt{f'} par \texttt{0,02f(15 − f)} il reste
\texttt{−0,02(15 − f)/f = −0,3/f + 0,02}, soit \texttt{−0,3g + 0,02}. C'est le seul endroit où
l'hypothèse \og{} \texttt{f} ne s'annule pas \fg{} sert.
\end{reponse}
\begin{proof}
La dérivée d'un quotient de la fonction constante $1$ par $F$ est
\[ \left(\frac{1}{F}\right)'(t)
   = \frac{0 \times F(t) - 1 \times F'(t)}{F(t)^2} = -\frac{F'(t)}{F(t)^2} , \]
ce qui suppose $F(t) \ne 0$ : c'est le seul endroit où cette hypothèse sert.

En y remplaçant $F'(t)$ par $0{,}02\,F(t)\,(15 - F(t))$, le facteur $F(t)$ se simplifie :
\[ -\frac{0{,}02\,F(t)\,(15 - F(t))}{F(t)^2}
   = -\frac{0{,}02\,(15 - F(t))}{F(t)}
   = -\frac{0{,}3}{F(t)} + 0{,}02 . \]

C'est bien $-0{,}3\,g(t) + 0{,}02$ avec $g = \dfrac{1}{F}$ : le changement de fonction a
rendu l'équation linéaire.
\end{proof}
\source{https://github.com/Commutator-IO/learn-lean/blob/main/exams/2025/Bac2025MetropoleJour1.lean\#L295}{Bac2025MetropoleJour1.lean\#L295}

\begin{reponse}
Partie B, question 3 : \texttt{f(t) = 15/(14e\textasciicircum{}(−0,3t) + 1)}.

Les solutions de \texttt{(E\(_{2}\))} sont les \texttt{t \(\mapsto\) Ce\textasciicircum{}(−0,3t) + 1/15}, et \texttt{g(0) = 1/f(0) = 1}
donne \texttt{C = 14/15}. Il ne reste qu'à inverser : \texttt{f = 1/g}, et le dénominateur ne
s'annule pas puisque l'exponentielle est positive.
\end{reponse}
\begin{proof}
Le dénominateur $14\mathrm{e}^{-0{,}3t} + 1$ est strictement positif, l'exponentielle
l'étant toujours ; on peut donc écrire l'égalité sous forme de produit et multiplier haut
et bas par $15$ :
\[ \frac{1}{\frac{14}{15}\mathrm{e}^{-0{,}3t} + \frac{1}{15}}
   = \frac{15}{14\mathrm{e}^{-0{,}3t} + 1} . \]
C'est la même fraction, écrite sans dénominateur au dénominateur.
\end{proof}
\source{https://github.com/Commutator-IO/learn-lean/blob/main/exams/2025/Bac2025MetropoleJour1.lean\#L310}{Bac2025MetropoleJour1.lean\#L310}

\begin{reponse}
Partie B, question 4 : la superficie tend vers \texttt{15 hectares}.

L'exponentielle \texttt{e\textasciicircum{}(−0,3t)} tend vers \texttt{0}, donc le dénominateur vers \texttt{1}. Le modèle
continu et le modèle discret donnent la même limite, ce qui est le seul point commun
qu'on leur demande.
\end{reponse}
\begin{proof}
La fonction $t \mapsto 0{,}3t$ tend vers $+\infty$, et l'exponentielle tend vers $0$ en
$-\infty$ : par composition, $\mathrm{e}^{-0{,}3t}$ tend vers $0$.

Donc le dénominateur $14\mathrm{e}^{-0{,}3t} + 1$ tend vers $14 \times 0 + 1 = 1$. Comme
cette limite n'est pas nulle, son inverse tend vers $\dfrac{1}{1} = 1$, et le produit par
$15$ tend vers $15$.

Le modèle continu donne donc la même limite que le modèle discret, $15$ hectares.
\end{proof}
\source{https://github.com/Commutator-IO/learn-lean/blob/main/exams/2025/Bac2025MetropoleJour1.lean\#L321}{Bac2025MetropoleJour1.lean\#L321}

\begin{reponse}
Partie B, question 5 : la superficie dépasse \texttt{14 hectares} après
\texttt{t > (10/3) ln 196}, soit environ \texttt{17,6} années \textemdash{} donc dans le courant de l'année
2042.

L'inéquation \texttt{15/(14e\textasciicircum{}(−0,3t) + 1) > 14} équivaut à \texttt{e\textasciicircum{}(−0,3t) < 1/196}, c'est-à-dire
\texttt{−0,3t < −ln 196}. Le seuil exact est \texttt{(ln 196)/0,3} : la valeur décimale ne se lit
qu'à la calculatrice, mais l'équivalence est exacte.
\end{reponse}
\begin{proof}
Le dénominateur $14\mathrm{e}^{-0{,}3t} + 1$ est strictement positif, donc l'inéquation
\[ \frac{15}{14\mathrm{e}^{-0{,}3t} + 1} > 14 \]
équivaut à $15 > 14\,(14\mathrm{e}^{-0{,}3t} + 1)$, c'est-à-dire à
$196\,\mathrm{e}^{-0{,}3t} < 1$, ou encore
\[ \mathrm{e}^{-0{,}3t} < \frac{1}{196} . \]

Le logarithme est strictement croissant et réciproque de l'exponentielle : cette dernière
inégalité équivaut à
\[ -0{,}3\,t < \ln\frac{1}{196} = -\ln 196 , \]
puis, en divisant par $0{,}3 > 0$, à $\dfrac{\ln 196}{0{,}3} < t$.

Le seuil exact est donc $\dfrac{\ln 196}{0{,}3} = \dfrac{10}{3}\ln 196 \approx 17{,}6$
années : la valeur décimale se lit à la calculatrice, l'équivalence est exacte.
\end{proof}
\source{https://github.com/Commutator-IO/learn-lean/blob/main/exams/2025/Bac2025MetropoleJour1.lean\#L340}{Bac2025MetropoleJour1.lean\#L340}

\section{Bac2025MetropoleJour2}

\noindent
Baccalauréat, épreuve de spécialité \textemdash{} voie générale, Métropole, 18 juin 2025 (jour 2).

Chaque question du sujet qui porte une proposition mathématique devient un théorème
Lean, dont la réponse est rédigée ci-dessous en français. Celles qui n'en portent pas \textemdash{} représenter un arbre pondéré, lire une durée
ou une asymptote sur un graphique, justifier une modélisation \textemdash{} n'ont pas d'énoncé, et
l'index du dossier dit lesquelles et pourquoi.

Les probabilités sont traitées comme le sujet les traite : sur les nombres. La
géométrie de l'espace, elle, est traitée sur les coordonnées : c'est ce que le repère
orthonormé permet, et le produit scalaire y est la somme des produits des coordonnées.

\subsection{Exercice 1 \textemdash{} chutes au roller, temps d'attente}

\begin{reponse}
Partie A, question 2 : \texttt{P(Ā \(\cap\) \(\overline{B}\)) = 0,4 \(\times\) 0,6 = 0,24}.

C'est la lecture d'un chemin de l'arbre : \texttt{24 \%} des personnes ne chutent à aucune des
deux séances.
\end{reponse}
\begin{proof}
Les probabilités se multiplient le long d'un chemin de l'arbre :
\[ P(\bar{A} \cap \bar{B}) = 0{,}4 \times 0{,}6 = 0{,}24 . \]
\end{proof}
\source{https://github.com/Commutator-IO/learn-lean/blob/main/exams/2025/Bac2025MetropoleJour2.lean\#L24}{Bac2025MetropoleJour2.lean\#L24}

\begin{reponse}
Partie A, question 3 : \texttt{P(B) = 0,34}.

La formule des probabilités totales décompose \texttt{B} selon qu'il y a eu chute ou non à la
première séance : \texttt{0,6 \(\times\) 0,3 + 0,4 \(\times\) 0,4}. Les deux chemins ne s'excluent pas par
hasard : \texttt{A} et \texttt{Ā} forment une partition.
\end{reponse}
\begin{proof}
Les évènements $A$ et $\bar{A}$ forment une partition : toute personne a chuté ou non à la
première séance, et pas les deux. La formule des probabilités totales décompose $B$ sur
cette partition :
\[ P(B) = P(A) \times P_A(B) + P(\bar{A}) \times P_{\bar{A}}(B)
   = 0{,}6 \times 0{,}3 + 0{,}4 \times 0{,}4 = 0{,}18 + 0{,}16 = 0{,}34 . \]
\end{proof}
\source{https://github.com/Commutator-IO/learn-lean/blob/main/exams/2025/Bac2025MetropoleJour2.lean\#L31}{Bac2025MetropoleJour2.lean\#L31}

\begin{reponse}
Partie A, question 4 : \texttt{P\_\(\overline{B}\)(Ā) = 0,24/0,66 = 4/11}, soit \texttt{0,364} au millième près.

La probabilité conditionnelle renverse le conditionnement : on connaît la probabilité
de chuter sachant qu'on a chuté avant, on cherche celle de n'avoir pas chuté avant
sachant qu'on ne chute pas maintenant.
\end{reponse}
\begin{proof}
Par définition d'une probabilité conditionnelle, et puisque
$P(\bar{B}) = 1 - P(B) = 0{,}66$ :
\[ P_{\bar{B}}(\bar{A}) = \frac{P(\bar{A} \cap \bar{B})}{P(\bar{B})}
   = \frac{0{,}24}{0{,}66} = \frac{24}{66} = \frac{4}{11} . \]

L'arrondi au millième se lit sur l'écart :
\[ \left|\frac{4}{11} - 0{,}364\right| = |0{,}363636\ldots - 0{,}364|
   = 0{,}000363\ldots < 0{,}0005 . \]

Le conditionnement est renversé : on connaît la probabilité de chuter sachant qu'on a
chuté, on cherche celle de n'avoir pas chuté sachant qu'on ne chute pas.
\end{proof}
\source{https://github.com/Commutator-IO/learn-lean/blob/main/exams/2025/Bac2025MetropoleJour2.lean\#L38}{Bac2025MetropoleJour2.lean\#L38}

\begin{reponse}
Partie A, question 5c : \texttt{E(X) = 100 \(\times\) 0,24 = 24}.

Sur un échantillon de \texttt{100} personnes, on en attend \texttt{24} qui ne chutent à aucune des
deux séances : l'espérance d'une loi binomiale est \texttt{np}, et c'est la proportion
retrouvée à l'échelle de l'échantillon.
\end{reponse}
\begin{proof}
L'espérance d'une loi binomiale de paramètres $n$ et $p$ vaut $np$ :
\[ E(X) = 100 \times 0{,}24 = 24 . \]
La proportion se retrouve à l'échelle de l'échantillon.
\end{proof}
\source{https://github.com/Commutator-IO/learn-lean/blob/main/exams/2025/Bac2025MetropoleJour2.lean\#L49}{Bac2025MetropoleJour2.lean\#L49}

\begin{reponse}
Partie B, question 1 : \texttt{E(T) = E(T\(_{1}\)) + E(T\(_{2}\)) = 100} minutes.

L'espérance est additive sans autre hypothèse : l'indépendance ne sert pas ici.
\end{reponse}
\begin{proof}
L'espérance est additive, sans aucune hypothèse sur les deux variables :
\[ E(T) = E(T_1) + E(T_2) = 40 + 60 = 100 . \]
\end{proof}
\source{https://github.com/Commutator-IO/learn-lean/blob/main/exams/2025/Bac2025MetropoleJour2.lean\#L54}{Bac2025MetropoleJour2.lean\#L54}

\begin{reponse}
Partie B, question 2 : \texttt{V(T) = 10\(^{2}\) + 16\(^{2}\) = 356}.

Les variances s'ajoutent parce que les deux variables sont indépendantes, et ce sont
les variances qui s'ajoutent, non les écarts-types : \texttt{10 + 16 = 26} n'a pas de sens
ici, et \texttt{\(\sqrt{356}\) \(\approx\) 18,9} est l'écart-type de la somme.
\end{reponse}
\begin{proof}
Les deux variables étant indépendantes, leurs variances s'ajoutent \textemdash{} et ce
sont les variances, non les écarts-types :
\[ V(T) = 10^2 + 16^2 = 100 + 256 = 356 . \]
L'écart-type de la somme est $\sqrt{356} \approx 18{,}9$, et non $10 + 16 = 26$.
\end{proof}
\source{https://github.com/Commutator-IO/learn-lean/blob/main/exams/2025/Bac2025MetropoleJour2.lean\#L61}{Bac2025MetropoleJour2.lean\#L61}

\begin{reponse}
Partie B, question 3 : \texttt{P(60 < T < 140) > 0,77}.

L'évènement \texttt{60 < T < 140} est exactement \texttt{|T − 100| < 40}. Bienaymé-Tchebychev majore
la probabilité contraire par \texttt{V(T)/40\(^{2}\) = 356/1 600 = 0,222 5}, donc minore celle
cherchée par \texttt{0,777 5}, qui dépasse \texttt{0,77}.

L'inégalité est très large \textemdash{} la vraie probabilité est sans doute bien plus grande \textemdash{}
mais elle ne demande aucune hypothèse sur la loi, et c'est là son intérêt.
\end{reponse}
\begin{proof}
L'évènement $60 < T < 140$ s'écrit $|T - 100| < 40$, puisque $E(T) = 100$. L'inégalité de
Bienaymé-Tchebychev majore la probabilité de l'évènement contraire :
\[ P(|T - 100| \ge 40) \le \frac{V(T)}{40^2} = \frac{356}{1\,600} = 0{,}2225 . \]

En passant au complémentaire,
\[ P(60 < T < 140) \ge 1 - 0{,}2225 = 0{,}7775 > 0{,}77 . \]

La majoration est très large \textemdash{} la vraie probabilité est sans doute bien plus
grande \textemdash{} mais elle ne suppose rien sur la loi de $T$, et c'est là son intérêt.
\end{proof}
\source{https://github.com/Commutator-IO/learn-lean/blob/main/exams/2025/Bac2025MetropoleJour2.lean\#L71}{Bac2025MetropoleJour2.lean\#L71}

\subsection{Exercice 2 \textemdash{} droites, plan et projeté orthogonal}

\begin{reponse}
Partie A, question 1 : les droites \texttt{d} et \texttt{d'} sont sécantes au point
\texttt{S(−1/2 ; 1 ; 4)}.

Le paramètre \texttt{t = −1} place le point de \texttt{d} en \texttt{S}, le paramètre \texttt{s = −1/2} place celui
de \texttt{d'} au même endroit : le point est commun. Les vecteurs directeurs \texttt{(2 ; 1 ; −1)} et
\texttt{(1 ; 1 ; −2)} n'étant pas colinéaires, les droites sont bien sécantes et non
confondues.
\end{reponse}
\begin{proof}
Il suffit d'exhiber les deux paramètres. Pour $t = -1$, le point de $d$ est
\[ \left(\frac{3}{2} + 2 \times (-1) ;\ 2 + (-1) ;\ 3 - (-1)\right)
   = \left(-\frac{1}{2} ;\ 1 ;\ 4\right) , \]
et pour $s = -\dfrac{1}{2}$, le point de $d'$ est
\[ \left(-\frac{1}{2} ;\ \frac{3}{2} - \frac{1}{2} ;\ 3 - 2 \times \left(-\frac{1}{2}\right)\right)
   = \left(-\frac{1}{2} ;\ 1 ;\ 4\right) . \]

Les deux droites passent donc par $S\left(-\tfrac{1}{2} ; 1 ; 4\right)$. Leurs vecteurs
directeurs $(2 ; 1 ; -1)$ et $(1 ; 1 ; -2)$ n'étant pas colinéaires, elles sont sécantes
et non confondues.
\end{proof}
\source{https://github.com/Commutator-IO/learn-lean/blob/main/exams/2025/Bac2025MetropoleJour2.lean\#L84}{Bac2025MetropoleJour2.lean\#L84}

\begin{reponse}
Partie A, question 2a : le vecteur \texttt{n(1 ; 2 ; 4)} est normal au plan \texttt{(ABC)}.

Il est orthogonal à \texttt{AB(2 ; −3 ; 1)} et à \texttt{AC(2 ; −1 ; 0)}, qui ne sont pas colinéaires :
deux orthogonalités à deux vecteurs non colinéaires du plan, c'est la définition d'un
vecteur normal.
\end{reponse}
\begin{proof}
Le produit scalaire est la somme des produits des coordonnées. Avec
$\vec{n}(1 ; 2 ; 4)$ :
\[ \vec{n} \cdot \vec{AB} = 1 \times 2 + 2 \times (-3) + 4 \times 1 = 2 - 6 + 4 = 0 , \]
\[ \vec{n} \cdot \vec{AC} = 1 \times 2 + 2 \times (-1) + 4 \times 0 = 2 - 2 + 0 = 0 . \]

Les vecteurs $\vec{AB}(2 ; -3 ; 1)$ et $\vec{AC}(2 ; -1 ; 0)$ ne sont pas colinéaires :
orthogonal à deux vecteurs non colinéaires du plan, $\vec{n}$ en est un vecteur normal.
\end{proof}
\source{https://github.com/Commutator-IO/learn-lean/blob/main/exams/2025/Bac2025MetropoleJour2.lean\#L94}{Bac2025MetropoleJour2.lean\#L94}

\begin{reponse}
Partie A, question 2b : une équation cartésienne du plan \texttt{(ABC)} est
\texttt{x + 2y + 4z − 7 = 0}.

Les coefficients sont les coordonnées du vecteur normal ; le terme constant se
détermine en écrivant qu'un point du plan vérifie l'équation. Les trois points la
vérifient, ce qui la confirme.
\end{reponse}
\begin{proof}
Un plan de vecteur normal $(1 ; 2 ; 4)$ a une équation de la forme $x + 2y + 4z + d = 0$ ;
le terme constant se détermine en écrivant qu'un point du plan la vérifie. Les trois
points la vérifient avec $d = -7$ :
\[ -1 + 2 \times 2 + 4 \times 1 - 7 = 0, \qquad
   1 + 2 \times (-1) + 4 \times 2 - 7 = 0, \]
\[ 1 + 2 \times 1 + 4 \times 1 - 7 = 0 . \]
\end{proof}
\source{https://github.com/Commutator-IO/learn-lean/blob/main/exams/2025/Bac2025MetropoleJour2.lean\#L104}{Bac2025MetropoleJour2.lean\#L104}

\begin{reponse}
Partie A, question 3 : les points \texttt{A}, \texttt{B}, \texttt{C} et \texttt{S} ne sont pas coplanaires.

Il suffit que \texttt{S} ne vérifie pas l'équation du plan \texttt{(ABC)} : la valeur \texttt{21/2} obtenue
est l'écart, non nul. Les quatre points forment donc un tétraèdre.
\end{reponse}
\begin{proof}
Il suffit de reporter les coordonnées de $S$ dans l'équation du plan :
\[ -\frac{1}{2} + 2 \times 1 + 4 \times 4 - 7 = \frac{21}{2} \ne 0 . \]
Le point $S$ n'appartient donc pas au plan $(ABC)$, et les quatre points forment un
tétraèdre.
\end{proof}
\source{https://github.com/Commutator-IO/learn-lean/blob/main/exams/2025/Bac2025MetropoleJour2.lean\#L113}{Bac2025MetropoleJour2.lean\#L113}

\begin{reponse}
Partie A, question 4a : \texttt{H(−1 ; 0 ; 2)} est le projeté orthogonal de \texttt{S} sur le
plan \texttt{(ABC)}.

Deux vérifications : \texttt{H} appartient au plan, et \texttt{SH(−1/2 ; −1 ; −2)} est colinéaire au
vecteur normal \textemdash{} c'est \texttt{−1/2 \(\times\) n}. Le projeté orthogonal est le seul point du plan tel
que la droite \texttt{(SH)} soit orthogonale au plan.
\end{reponse}
\begin{proof}
Deux vérifications suffisent, car le projeté orthogonal est le seul point du plan tel que
la droite qui le joint à $S$ soit orthogonale au plan.

D'abord $H(-1 ; 0 ; 2)$ appartient au plan :
\[ -1 + 2 \times 0 + 4 \times 2 - 7 = 0 . \]

Ensuite $\vec{SH}$ est colinéaire au vecteur normal :
\[ \vec{SH} = \left(-1 + \frac{1}{2} ;\ 0 - 1 ;\ 2 - 4\right)
   = \left(-\frac{1}{2} ;\ -1 ;\ -2\right) = -\frac{1}{2}\,(1 ;\ 2 ;\ 4) . \]
\end{proof}
\source{https://github.com/Commutator-IO/learn-lean/blob/main/exams/2025/Bac2025MetropoleJour2.lean\#L121}{Bac2025MetropoleJour2.lean\#L121}

\begin{reponse}
Partie A, question 4b : aucun point \texttt{M} du plan \texttt{(ABC)} ne vérifie \texttt{SM < \(\sqrt{21}\)/2}.

La distance de \texttt{S} au plan est \texttt{SH = \(\sqrt{1/4 + 1 + 4}\) = \(\sqrt{21}\)/2}, et le projeté orthogonal
réalise le minimum. Plutôt que d'invoquer ce minimum, on le démontre : pour tout point
\texttt{M(x ; y ; z)} du plan, l'inégalité de Cauchy-Schwarz appliquée à
\texttt{(x + 1/2) + 2(y − 1) + 4(z − 4) = −21/2} donne \texttt{(21/2)\(^{2}\) \(\le\) 21 \(\times\) SM\(^{2}\)}, soit
\texttt{SM\(^{2}\) \(\ge\) 21/4}.

C'est le carré de la distance qui est comparé, jamais la distance elle-même : la
racine carrée n'apparaît pas dans la démonstration.
\end{reponse}

\begin{figure}[h]
\centering
% Construction (perspective cavalière) : le plan (ABC) est le parallélogramme
% (0 ; 0), (6 ; 0), (7,6 ; 2), (1,6 ; 2). H = (3,2 ; 0,9), M = (5,4 ; 1,5),
% et S = (3,2 ; 3,7), à la verticale de H sur le dessin.
% [SH] est perpendiculaire au plan : c'est la plus courte distance de S au plan.
\begin{tikzpicture}[scale=0.78]
  \draw (0, 0) -- (6, 0) -- (7.6, 2) -- (1.6, 2) -- cycle;
  \coordinate (H) at (3.2, 0.9);
  \coordinate (M) at (5.4, 1.5);
  \coordinate (S) at (3.2, 3.7);
  \draw[dashed] (S) -- (H);
  \draw (S) -- (M);
  \draw[dashed] (H) -- (M);
  \draw (3.2, 1.18) -- (3.48, 1.18) -- (3.48, 0.9);
  \foreach \p in {H, M, S} \fill (\p) circle (1.6pt);
  \node[above] at (S) {$S$};
  \node[below right] at (H) {$H$};
  \node[below right] at (M) {$M$};
  \node[left] at (3.2, 2.4) {$\dfrac{\sqrt{21}}{2}$};
  \node[right] at (7, 0.4) {$(ABC)$};
\end{tikzpicture}

\caption{Le projeté orthogonal \(H\) réalise la plus courte distance de \(S\) au plan : tout autre point \(M\) est plus loin.}
\end{figure}
\begin{proof}
Posons $a = x + \dfrac{1}{2}$, $b = y - 1$, $c = z - 4$, de sorte que $SM^2 = a^2 + b^2 +
c^2$. L'hypothèse $x + 2y + 4z = 7$ se récrit
\[ a + 2b + 4c = (x + 2y + 4z) - \frac{35}{2} = 7 - \frac{35}{2} = -\frac{21}{2} . \]

L'identité de Lagrange, pour les triplets $(1 ; 2 ; 4)$ et $(a ; b ; c)$, s'écrit
\[ (1 + 4 + 16)(a^2 + b^2 + c^2) - (a + 2b + 4c)^2
   = (b - 2a)^2 + (c - 4a)^2 + (2c - 4b)^2 \ge 0 , \]
somme de trois carrés. C'est l'inégalité de Cauchy-Schwarz, démontrée plutôt qu'invoquée.

D'où $\left(\dfrac{21}{2}\right)^2 \le 21\,(a^2 + b^2 + c^2)$, c'est-à-dire
\[ SM^2 = a^2 + b^2 + c^2 \ge \frac{441}{4 \times 21} = \frac{21}{4} . \]

Aucun point du plan n'est donc à une distance de $S$ inférieure à
$\sqrt{21}/2$, distance réalisée par le projeté $H$. C'est le carré de la distance qui est
comparé d'un bout à l'autre : la racine carrée n'intervient jamais.
\end{proof}
\source{https://github.com/Commutator-IO/learn-lean/blob/main/exams/2025/Bac2025MetropoleJour2.lean\#L136}{Bac2025MetropoleJour2.lean\#L136}

\begin{reponse}
Partie B, question 1 : le point \texttt{M} de \texttt{[CS]} tel que \texttt{CM = k CS} a pour
coordonnées \texttt{(1 − 3k/2 ; 1 ; 1 + 3k)}.

L'ordonnée ne dépend pas de \texttt{k} : \texttt{C} et \texttt{S} ont la même, et le segment est donc
contenu dans le plan \texttt{y = 1}.
\end{reponse}
\begin{proof}
Les coordonnées de $M$ s'obtiennent en ajoutant à celles de $C$ le vecteur
$k\,\vec{CS}$, et il n'y a plus qu'à développer :
\[ 1 + k\left(-\frac{1}{2} - 1\right) = 1 - \frac{3}{2}k , \qquad
   1 + k\,(1 - 1) = 1 , \qquad
   1 + k\,(4 - 1) = 1 + 3k . \]

L'ordonnée ne dépend pas de $k$ : $C$ et $S$ ont la même, donc le segment est contenu dans
le plan d'équation $y = 1$.
\end{proof}
\source{https://github.com/Commutator-IO/learn-lean/blob/main/exams/2025/Bac2025MetropoleJour2.lean\#L146}{Bac2025MetropoleJour2.lean\#L146}

\begin{reponse}
Partie B, question 2 : oui, il existe un point \texttt{M} de \texttt{[CS]} tel que le triangle
\texttt{MAB} soit rectangle en \texttt{M}.

Le produit scalaire \texttt{MA \(\cdot\) MB} vaut \texttt{11,25k\(^{2}\) − 6k − 2}. Il s'annule pour
\texttt{k = (12 + 2\(\sqrt{126}\))/45 \(\approx\) 0,766}, qui appartient bien à \texttt{[0 ; 1]} \textemdash{} l'autre racine est
négative, donc hors du segment. Le point cherché existe et il est unique.

La racine est irrationnelle : le sujet demande une existence, non une valeur, et c'est
tout l'intérêt de la question.
\end{reponse}
\begin{proof}
En développant, le produit scalaire $\vec{MA} \cdot \vec{MB}$ vaut
\[ (-2 + 1{,}5k)(1{,}5k) + 1 \times (-2) + (-3k)(1 - 3k) = 11{,}25\,k^2 - 6k - 2 . \]

Le discriminant de ce trinôme vaut $36 + 90 = 126$, et ses racines sont
$\dfrac{6 \pm \sqrt{126}}{11{,}25 \times 2} = \dfrac{12 \pm 2\sqrt{126}}{45}$.

Prenons $k = \dfrac{12 + 2\sqrt{126}}{45}$ et vérifions qu'il appartient bien à
$[0 ; 1]$. D'une part $\sqrt{126} \ge 0$, donc $k \ge \dfrac{12}{45} > 0$. D'autre part
$126 \le 144$ donne $\sqrt{126} \le 12$, donc
\[ k \le \frac{12 + 24}{45} = \frac{36}{45} < 1 . \]

Il reste à voir que ce $k$ annule le trinôme, ce qui est un calcul où
$\sqrt{126}^{\,2} = 126$ est la seule chose à savoir.

Le point cherché existe donc, et il est unique : l'autre racine est négative, donc hors du
segment. Le sujet demande une existence, non une valeur, et c'est tout l'intérêt de la
question \textemdash{} la racine est irrationnelle.
\end{proof}
\source{https://github.com/Commutator-IO/learn-lean/blob/main/exams/2025/Bac2025MetropoleJour2.lean\#L159}{Bac2025MetropoleJour2.lean\#L159}

\subsection{Exercice 3 \textemdash{} vrai ou faux : suites, convexité, logarithme}

\begin{reponse}
Affirmation 1, \textbf{vraie} : la suite \texttt{u\_n = (1 + 5n)/(2 + 3n)} converge vers \texttt{5/3}.

La forme \texttt{5/3 − 7/(9n + 6)} rend la limite immédiate : l'écart à \texttt{5/3} est un quotient
de constante par une quantité qui tend vers \texttt{+\(\infty\)}. C'est le rapport des coefficients
dominants, et l'écriture le démontre au lieu de l'invoquer.
\end{reponse}
\begin{proof}
On récrit le terme général en isolant l'écart à la limite. Pour tout entier $n$, le
dénominateur $2 + 3n$ est strictement positif, et
\[ \frac{5}{3} - \frac{7}{9n + 6}
  = \frac{5(9n + 6) - 21}{3\,(9n + 6)}
  = \frac{45n + 9}{27n + 18}
  = \frac{9\,(5n + 1)}{9\,(3n + 2)}
  = \frac{1 + 5n}{2 + 3n} . \]

Or $9n + 6$ tend vers $+\infty$, donc $\dfrac{7}{9n + 6}$ tend vers $0$ : le quotient
d'une constante par une quantité qui tend vers $+\infty$ tend vers $0$.

Par différence, $u_n$ tend vers $\dfrac{5}{3} - 0 = \dfrac{5}{3}$.

C'est le rapport des coefficients dominants, mais l'écriture le démontre au lieu de
l'invoquer.
\end{proof}
\source{https://github.com/Commutator-IO/learn-lean/blob/main/exams/2025/Bac2025MetropoleJour2.lean\#L176}{Bac2025MetropoleJour2.lean\#L176}

\begin{definition}
La suite de l'affirmation 2 : \texttt{w\(_{0}\) = 0} et \texttt{w\_\{n+1\} = 3w\_n − 2n + 3}.
\end{definition}
\source{https://github.com/Commutator-IO/learn-lean/blob/main/exams/2025/Bac2025MetropoleJour2.lean\#L197}{Bac2025MetropoleJour2.lean\#L197}

\begin{reponse}
La relation de récurrence, sous la forme où la preuve l'utilise.
\end{reponse}
\begin{proof}
Les deux membres sont le même énoncé : c'est la seconde clause de la définition récursive
de $w$.
\end{proof}
\source{https://github.com/Commutator-IO/learn-lean/blob/main/exams/2025/Bac2025MetropoleJour2.lean\#L202}{Bac2025MetropoleJour2.lean\#L202}

\begin{reponse}
Affirmation 2, \textbf{vraie} : pour tout entier naturel \texttt{n}, \texttt{w\_n \(\ge\) n}.

La récurrence est immédiate et instructive : de \texttt{w\_n \(\ge\) n} on tire
\texttt{w\_\{n+1\} \(\ge\) 3n − 2n + 3 = n + 3}, qui dépasse \texttt{n + 1} de deux unités. La marge grandit,
ce qui est le signe qu'une minoration bien plus forte serait vraie \textemdash{} mais celle-ci
suffit à répondre.
\end{reponse}
\begin{proof}
Par récurrence sur $n$.

\emph{Initialisation.} $w_0 = 0$, donc $0 \le w_0$.

\emph{Hérédité.} Supposons $n \le w_n$. Alors
\[ w_{n+1} = 3w_n - 2n + 3 \ge 3n - 2n + 3 = n + 3 \ge n + 1 . \]

La marge grandit de deux unités à chaque rang, signe qu'une minoration bien plus forte
serait vraie ; celle-ci suffit à répondre.
\end{proof}
\source{https://github.com/Commutator-IO/learn-lean/blob/main/exams/2025/Bac2025MetropoleJour2.lean\#L210}{Bac2025MetropoleJour2.lean\#L210}

\begin{reponse}
Affirmation 4, \textbf{vraie} : pour tout réel \texttt{x > 0}, \texttt{ln x − x + 1 \(\le\) 0}.

C'est l'inégalité de concavité du logarithme, \texttt{ln x \(\le\) x − 1}, avec égalité en \texttt{1}
seulement : la courbe du logarithme est sous sa tangente en \texttt{1}.
\end{reponse}

\begin{figure}[h]
\centering
% Construction : y = ln x sur ]0 ; 3,6], et la droite y = x − 1, tangente à la
% courbe au point (1 ; 0). La courbe reste sous la droite, avec égalité en ce
% seul point : c'est l'inégalité ln x \(\le\) x − 1.
\begin{tikzpicture}[xscale=1.15, yscale=0.8]
  \draw[->] (0, 0) -- (3.9, 0) node[below] {$x$};
  \draw[->] (0, -2.1) -- (0, 2.8);
  \draw plot coordinates {(0.140, -1.966) (0.198, -1.621) (0.255, -1.365) (0.313, -1.162) (0.371, -0.992) (0.428, -0.848) (0.486, -0.722) (0.544, -0.609) (0.601, -0.509) (0.659, -0.417) (0.717, -0.333) (0.774, -0.256) (0.832, -0.184) (0.890, -0.117) (0.947, -0.054) (1.005, 0.005) (1.063, 0.061) (1.120, 0.114) (1.178, 0.164) (1.236, 0.212) (1.293, 0.257) (1.351, 0.301) (1.409, 0.343) (1.466, 0.383) (1.524, 0.421) (1.582, 0.458) (1.639, 0.494) (1.697, 0.529) (1.755, 0.562) (1.812, 0.595) (1.870, 0.626) (1.928, 0.656) (1.985, 0.686) (2.043, 0.714) (2.101, 0.742) (2.158, 0.769) (2.216, 0.796) (2.274, 0.821) (2.331, 0.846) (2.389, 0.871) (2.447, 0.895) (2.504, 0.918) (2.562, 0.941) (2.620, 0.963) (2.677, 0.985) (2.735, 1.006) (2.793, 1.027) (2.850, 1.047) (2.908, 1.067) (2.966, 1.087) (3.023, 1.106) (3.081, 1.125) (3.139, 1.144) (3.196, 1.162) (3.254, 1.180) (3.312, 1.197) (3.369, 1.215) (3.427, 1.232) (3.485, 1.248) (3.542, 1.265) (3.600, 1.281)};
  \draw[dashed] (-0.6, -1.6) -- (3.6, 2.6);
  \fill (1, 0) circle (1.1pt);
  \node[below right] at (1, 0) {$1$};
  \node[right] at (3.5, 1.28) {$y = \ln x$};
  \node[right] at (3.35, 2.6) {$y = x - 1$};
\end{tikzpicture}

\caption{La courbe du logarithme reste sous sa tangente au point d'abscisse \(1\).}
\end{figure}
\begin{proof}
C'est l'inégalité $\ln x \le x - 1$, valable pour tout $x > 0$, qu'il reste à réorganiser.

Cette inégalité s'obtient elle-même de la convexité de l'exponentielle, sous la forme
$1 + u \le \mathrm{e}^u$ : la fonction $u \mapsto \mathrm{e}^u - u - 1$ a pour dérivée
$\mathrm{e}^u - 1$, qui est négative avant $0$ et positive après ; son minimum est donc
atteint en $0$ et vaut $0$. En appliquant cela à $u = \ln x$, licite puisque $x > 0$ :
\[ 1 + \ln x \le \mathrm{e}^{\ln x} = x , \]
c'est-à-dire $\ln x - x + 1 \le 0$.

Géométriquement, la courbe du logarithme est sous sa tangente en $1$, avec égalité en ce
seul point.
\end{proof}
\source{https://github.com/Commutator-IO/learn-lean/blob/main/exams/2025/Bac2025MetropoleJour2.lean\#L222}{Bac2025MetropoleJour2.lean\#L222}

\subsection{Exercice 4 \textemdash{} le freinage du chariot}

\begin{reponse}
Partie B, question 1b : \texttt{g(t) = te\textasciicircum{}(−0,6t)} est solution de
\texttt{(E) : y' + 0,6y = e\textasciicircum{}(−0,6t)}.

La dérivée d'un produit donne \texttt{g'(t) = e\textasciicircum{}(−0,6t) − 0,6te\textasciicircum{}(−0,6t)}, c'est-à-dire
\texttt{e\textasciicircum{}(−0,6t) − 0,6g(t)} : l'équation est vérifiée telle quelle.
\end{reponse}
\begin{proof}
La fonction $s \mapsto -0{,}6s$ a pour dérivée $-0{,}6$, donc $s \mapsto
\mathrm{e}^{-0{,}6s}$ a pour dérivée $-0{,}6\,\mathrm{e}^{-0{,}6s}$.

La dérivée d'un produit donne alors
\[ g'(t) = 1 \times \mathrm{e}^{-0{,}6t} + t \times (-0{,}6)\,\mathrm{e}^{-0{,}6t}
   = \mathrm{e}^{-0{,}6t} - 0{,}6\,t\,\mathrm{e}^{-0{,}6t} , \]
c'est-à-dire $g'(t) = \mathrm{e}^{-0{,}6t} - 0{,}6\,g(t)$ : l'équation
$y' + 0{,}6y = \mathrm{e}^{-0{,}6t}$ est vérifiée telle quelle.
\end{proof}
\source{https://github.com/Commutator-IO/learn-lean/blob/main/exams/2025/Bac2025MetropoleJour2.lean\#L233}{Bac2025MetropoleJour2.lean\#L233}

\begin{reponse}
Partie B, question 2a : \texttt{v'(t) = (−6,2 − 0,6t)e\textasciicircum{}(−0,6t)}, où
\texttt{v(t) = (12 + t)e\textasciicircum{}(−0,6t)}.

C'est encore la dérivée d'un produit : \texttt{e\textasciicircum{}(−0,6t) − 0,6(12 + t)e\textasciicircum{}(−0,6t)}, et
\texttt{1 − 7,2 = −6,2}.
\end{reponse}
\begin{proof}
Même dérivation d'un produit, avec le facteur $12 + s$ dont la dérivée vaut $1$ :
\[ v'(t) = \mathrm{e}^{-0{,}6t} + (12 + t) \times (-0{,}6)\,\mathrm{e}^{-0{,}6t}
   = \bigl(1 - 7{,}2 - 0{,}6t\bigr)\mathrm{e}^{-0{,}6t}
   = (-6{,}2 - 0{,}6t)\,\mathrm{e}^{-0{,}6t} . \]
\end{proof}
\source{https://github.com/Commutator-IO/learn-lean/blob/main/exams/2025/Bac2025MetropoleJour2.lean\#L247}{Bac2025MetropoleJour2.lean\#L247}

\begin{reponse}
Partie B, question 1d : \texttt{v(t) = (12 + t)e\textasciicircum{}(−0,6t)} est bien la solution cherchée \textemdash{}
elle vérifie l'équation et la condition initiale \texttt{v(0) = 12}.

La vérification de l'équation est algébrique : \texttt{−6,2 − 0,6t + 0,6(12 + t) = 1}. La
condition initiale est ce qui détermine la constante, et donc ce qui rend la solution
unique.
\end{reponse}
\begin{proof}
La vérification de l'équation est algébrique : l'exponentielle se met en facteur et il
reste
\[ (-6{,}2 - 0{,}6t) + 0{,}6\,(12 + t) = -6{,}2 - 0{,}6t + 7{,}2 + 0{,}6t = 1 , \]
de sorte que $v'(t) + 0{,}6\,v(t) = \mathrm{e}^{-0{,}6t}$.

La condition initiale se lit directement, $\mathrm{e}^0 = 1$ :
\[ v(0) = (12 + 0)\,\mathrm{e}^{0} = 12 . \]
C'est elle qui détermine la constante, donc qui rend la solution unique.
\end{proof}
\source{https://github.com/Commutator-IO/learn-lean/blob/main/exams/2025/Bac2025MetropoleJour2.lean\#L264}{Bac2025MetropoleJour2.lean\#L264}

\begin{reponse}
Partie B, question 2c : la vitesse est strictement décroissante sur \texttt{[0 ; +\(\infty\)[}.

La dérivée est le produit de \texttt{−6,2 − 0,6t}, négatif dès que \texttt{t \(\ge\) 0}, par une
exponentielle, toujours positive. Le chariot ralentit sans jamais s'arrêter : c'est
pourquoi un dispositif mécanique est nécessaire.
\end{reponse}
\begin{proof}
Pour $t \ge 0$, le facteur $-6{,}2 - 0{,}6t$ est strictement négatif, puisque
$0{,}6t \ge 0$. L'exponentielle, elle, est toujours strictement positive. Le produit d'un
nombre strictement négatif par un nombre strictement positif est strictement négatif :
$v'(t) < 0$.

La vitesse décroît donc strictement, sans jamais s'annuler : c'est pourquoi un dispositif
mécanique d'arrêt est nécessaire.
\end{proof}
\source{https://github.com/Commutator-IO/learn-lean/blob/main/exams/2025/Bac2025MetropoleJour2.lean\#L275}{Bac2025MetropoleJour2.lean\#L275}

\begin{reponse}
Partie B, question 2b : la vitesse tend vers \texttt{0} en \texttt{+\(\infty\)}.

L'exponentielle l'emporte sur le terme linéaire, ce que le sujet fait admettre en
récrivant \texttt{te\textasciicircum{}(−0,6t)} sous la forme \texttt{(1/0,6)(0,6t)/e\textasciicircum{}(0,6t)}. La croissance comparée
est ici invoquée sous sa forme mathlib : \texttt{x e\textasciicircum{}(−x)} tend vers \texttt{0}.
\end{reponse}
\begin{proof}
On sépare les deux termes :
\[ (12 + t)\,\mathrm{e}^{-0{,}6t}
   = 12\,\mathrm{e}^{-0{,}6t} + \frac{1}{0{,}6}\,(0{,}6t)\,\mathrm{e}^{-0{,}6t} . \]

Le premier tend vers $0$, l'exponentielle tendant vers $0$ en $-\infty$ et $0{,}6t$ vers
$+\infty$.

Le second aussi, par croissance comparée : $x\,\mathrm{e}^{-x}$ tend vers $0$ en
$+\infty$, l'exponentielle l'emportant sur le terme linéaire. Ce résultat est emprunté
\textemdash{} il repose sur la définition analytique de l'exponentielle et sa minoration
par $\dfrac{x^2}{2}$ \textemdash{} et le sujet le fait admettre de même, en récrivant
$t\,\mathrm{e}^{-0{,}6t}$ sous la forme $\dfrac{1}{0{,}6} \dfrac{0{,}6t}{\mathrm{e}^{0{,}6t}}$.

La somme tend donc vers $0$.
\end{proof}
\source{https://github.com/Commutator-IO/learn-lean/blob/main/exams/2025/Bac2025MetropoleJour2.lean\#L284}{Bac2025MetropoleJour2.lean\#L284}

\begin{reponse}
Partie C, question 1 : la distance parcourue entre les instants \texttt{0} et \texttt{t} vaut
\texttt{e\textasciicircum{}(−0,6t)(−5t/3 − 205/9) + 205/9}.

Le sujet demande une intégration par parties ; ce qui se vérifie ici est le résultat,
en dérivant la primitive proposée : sa dérivée est bien \texttt{(12 + t)e\textasciicircum{}(−0,6t)}, et elle
s'annule en \texttt{0}. Le théorème fondamental de l'analyse fait le reste.

La limite de cette expression en \texttt{+\(\infty\)} est \texttt{205/9 \(\approx\) 22,8} : c'est la longueur minimale
de la zone de freinage, celle que la partie A faisait lire sur l'asymptote.
\end{reponse}

\begin{figure}[h]
\centering
% Construction : v(t) = (12 + t)e^{−0,6t} sur [0 ; 8], ordonnées divisées par 3.
% L'aire hachurée, de 0 à t, est la distance parcourue pendant le freinage ;
% elle croît vers 205/9 \(\approx\) 22,8 mètres, longueur minimale de la zone d'arrêt.
\begin{tikzpicture}[xscale=1.05, yscale=0.85]
  \draw[->] (0, 0) -- (8.6, 0);
  \draw[->] (0, 0) -- (0, 4.6) node[left] {$v$};
  \draw[thin] (0.300, 0) -- (0.300, 3.425);
  \draw[thin] (0.600, 0) -- (0.600, 2.930);
  \draw[thin] (0.900, 0) -- (0.900, 2.506);
  \draw[thin] (1.200, 0) -- (1.200, 2.142);
  \draw[thin] (1.500, 0) -- (1.500, 1.830);
  \draw[thin] (1.800, 0) -- (1.800, 1.562);
  \draw[thin] (2.100, 0) -- (2.100, 1.333);
  \draw[thin] (2.400, 0) -- (2.400, 1.137);
  \draw[thin] (2.700, 0) -- (2.700, 0.970);
  \draw[thin] (3.000, 0) -- (3.000, 0.826);
  \draw[thin] (3.300, 0) -- (3.300, 0.704);
  \draw[thin] (3.600, 0) -- (3.600, 0.600);
  \draw[thin] (3.900, 0) -- (3.900, 0.511);
  \draw[thin] (4.200, 0) -- (4.200, 0.434);
  \draw[thick] plot coordinates {(0.000, 4.000) (0.167, 3.670) (0.333, 3.366) (0.500, 3.087) (0.667, 2.830) (0.833, 2.595) (1.000, 2.378) (1.167, 2.179) (1.333, 1.997) (1.500, 1.830) (1.667, 1.676) (1.833, 1.535) (2.000, 1.406) (2.167, 1.287) (2.333, 1.178) (2.500, 1.078) (2.667, 0.987) (2.833, 0.903) (3.000, 0.826) (3.167, 0.756) (3.333, 0.692) (3.500, 0.633) (3.667, 0.579) (3.833, 0.529) (4.000, 0.484) (4.167, 0.442) (4.333, 0.404) (4.500, 0.370) (4.667, 0.338) (4.833, 0.309) (5.000, 0.282) (5.167, 0.258) (5.333, 0.236) (5.500, 0.215) (5.667, 0.197) (5.833, 0.180) (6.000, 0.164) (6.167, 0.150) (6.333, 0.137) (6.500, 0.125) (6.667, 0.114) (6.833, 0.104) (7.000, 0.095) (7.167, 0.087) (7.333, 0.079) (7.500, 0.072) (7.667, 0.066) (7.833, 0.060) (8.000, 0.055)};
  \draw[dotted] (4.5, 0) -- (4.5, {(12 + 4.5) * exp(-2.7) / 3});
  \node[below] at (4.5, -0.05) {$t$};
  \node[right] at (5.4, 1.1) {$\mathcal{C}_v$};
\end{tikzpicture}

\caption{L'aire sous la courbe de la vitesse, de \(0\) à \(t\), est la distance parcourue pendant le freinage.}
\end{figure}
\begin{proof}
Posons $F(s) = \mathrm{e}^{-0{,}6s}\left(-\dfrac{5}{3}s - \dfrac{205}{9}\right) +
\dfrac{205}{9}$ et vérifions que $F$ est une primitive de la vitesse.

La dérivée d'un produit donne
\[ F'(s) = -0{,}6\,\mathrm{e}^{-0{,}6s}\left(-\frac{5}{3}s - \frac{205}{9}\right)
        + \mathrm{e}^{-0{,}6s}\left(-\frac{5}{3}\right)
   = \mathrm{e}^{-0{,}6s}\left(s + \frac{41}{3} - \frac{5}{3}\right)
   = (12 + s)\,\mathrm{e}^{-0{,}6s} , \]
puisque $0{,}6 \times \dfrac{5}{3} = 1$ et $0{,}6 \times \dfrac{205}{9} = \dfrac{41}{3}$.

La fonction intégrée est continue, donc intégrable sur le segment. Le théorème fondamental
de l'analyse donne alors
\[ \int_0^t (12 + x)\,\mathrm{e}^{-0{,}6x}\,\mathrm{d}x = F(t) - F(0)
   = \mathrm{e}^{-0{,}6t}\left(-\frac{5}{3}t - \frac{205}{9}\right) + \frac{205}{9} , \]
car $F(0) = -\dfrac{205}{9} + \dfrac{205}{9} = 0$.

Le sujet demandait une intégration par parties ; ce qui est vérifié ici est son résultat.
La limite de cette expression en $+\infty$ vaut $\dfrac{205}{9} \approx 22{,}8$ : c'est la
longueur minimale de la zone de freinage.
\end{proof}
\source{https://github.com/Commutator-IO/learn-lean/blob/main/exams/2025/Bac2025MetropoleJour2.lean\#L310}{Bac2025MetropoleJour2.lean\#L310}

\section{Brevet2025Metropole}

\noindent
Brevet des collèges \textemdash{} série générale, Métropole\textendash{}Antilles\textendash{}Guyane, 26 juin 2025.

Chaque question du sujet qui porte une proposition mathématique devient un théorème
Lean, dont la réponse est rédigée ci-dessous en français. Celles qui n'en portent pas \textemdash{} reconnaître une transformation sur une
figure, compléter un programme Scratch, écrire une formule de tableur \textemdash{} n'ont pas
d'énoncé, et l'index du dossier dit lesquelles et pourquoi.

Les énoncés sont ceux du sujet, pas leur généralisation. Quand la configuration
géométrique est donnée par une figure, ce qu'elle établit \textemdash{} l'angle droit, les
alignements \textemdash{} devient une hypothèse : la figure n'est pas formalisée, l'arithmétique
qui en découle l'est.

\subsection{Exercice 1 \textemdash{} deux urnes}

\begin{definition}
L'urne \texttt{A} contient les boules \texttt{7, 10, 12, 15, 24, 30}.
\end{definition}
\source{https://github.com/Commutator-IO/learn-lean/blob/main/exams/2025/Brevet2025Metropole.lean\#L22}{Brevet2025Metropole.lean\#L22}

\begin{definition}
L'urne \texttt{B} contient les boules \texttt{2, 5, 6, 8, 17, 18, 21, 22, 25}.
\end{definition}
\source{https://github.com/Commutator-IO/learn-lean/blob/main/exams/2025/Brevet2025Metropole.lean\#L25}{Brevet2025Metropole.lean\#L25}

\begin{reponse}
Question 1 : quatre des six boules de l'urne \texttt{A} portent un nombre pair, donc
la probabilité d'en tirer une vaut \texttt{2/3}.

Les boules étant indiscernables au toucher, la probabilité est le quotient du
nombre de cas favorables par le nombre de cas possibles : c'est l'équiprobabilité,
et c'est la seule hypothèse de tout l'exercice.
\end{reponse}
\begin{proof}
On parcourt les six boules de l'urne $A$ et l'on retient celles dont le reste dans la
division par $2$ est nul : $10$, $12$, $24$ et $30$. Il y en a quatre.

Les boules étant indiscernables au toucher, chacune a la même chance d'être tirée, et la
probabilité cherchée est le quotient du nombre de cas favorables par le nombre de cas
possibles :
\[ \frac{4}{6} = \frac{2 \times 2}{2 \times 3} = \frac{2}{3} . \]
\end{proof}
\source{https://github.com/Commutator-IO/learn-lean/blob/main/exams/2025/Brevet2025Metropole.lean\#L33}{Brevet2025Metropole.lean\#L33}

\begin{reponse}
Question 2 : trois des neuf boules de l'urne \texttt{B} portent un nombre premier \textemdash{}
\texttt{2}, \texttt{5} et \texttt{17} \textemdash{}, donc la probabilité vaut \texttt{1/3}.

Les six autres sont composées : \texttt{6}, \texttt{8}, \texttt{18} et \texttt{22} sont pairs et plus grands que
\texttt{2}, \texttt{21 = 3 \(\times\) 7}, et \texttt{25 = 5\(^{2}\)}.
\end{reponse}
\begin{proof}
Trois boules de $B$ portent un nombre premier. En effet $2$, $5$ et $17$ n'ont pas d'autre
diviseur que $1$ et eux-mêmes.

Les six autres ne sont pas premières : $6$, $8$, $18$ et $22$ sont pairs et plus grands que
$2$, donc divisibles par $2$ ; puis $21 = 3 \times 7$ et $25 = 5 \times 5$.

Il y a donc trois cas favorables sur neuf, et
\[ \frac{3}{9} = \frac{3 \times 1}{3 \times 3} = \frac{1}{3} . \]
\end{proof}
\source{https://github.com/Commutator-IO/learn-lean/blob/main/exams/2025/Brevet2025Metropole.lean\#L42}{Brevet2025Metropole.lean\#L42}

\begin{reponse}
Question 3 : l'urne \texttt{A} contient trois multiples de \texttt{6} \textemdash{} \texttt{12}, \texttt{24}, \texttt{30} \textemdash{} et
l'urne \texttt{B} deux seulement \textemdash{} \texttt{6} et \texttt{18}. C'est donc \texttt{A}.

La question porte sur un effectif, non sur une probabilité : c'est \texttt{A} qui en contient
le plus, alors que la proportion y est aussi la plus grande, \texttt{1/2} contre \texttt{2/9}.
\end{reponse}
\begin{proof}
On teste chaque boule : un nombre est multiple de $6$ quand le reste de sa division par $6$
est nul. Dans $A$, ce sont $12$, $24$ et $30$ \textemdash{} trois boules. Dans $B$, ce sont
$6$ et $18$ \textemdash{} deux boules.

C'est donc l'urne $A$ qui en contient le plus.
\end{proof}
\source{https://github.com/Commutator-IO/learn-lean/blob/main/exams/2025/Brevet2025Metropole.lean\#L55}{Brevet2025Metropole.lean\#L55}

\begin{reponse}
Question 4 : la probabilité d'obtenir un nombre supérieur ou égal à \texttt{20} est
\texttt{1/3} dans les deux urnes \textemdash{} \texttt{2} boules sur \texttt{6} d'un côté, \texttt{3} sur \texttt{9} de l'autre.

Les deux fractions ne sont pas les mêmes, mais elles ont la même valeur : c'est ce
que la question demande de démontrer, et rien n'oblige les urnes à avoir le même
nombre de boules pour cela.
\end{reponse}
\begin{proof}
Dans $A$, deux boules portent un nombre supérieur ou égal à $20$ : $24$ et $30$. Dans $B$,
il y en a trois : $21$, $22$ et $25$.

Les deux probabilités sont $\dfrac{2}{6}$ et $\dfrac{3}{9}$ : elles ne s'écrivent pas de la
même façon, mais les produits en croix sont égaux,
\[ 2 \times 9 = 18 = 6 \times 3 , \]
donc les deux fractions ont la même valeur, $\dfrac{1}{3}$.
\end{proof}
\source{https://github.com/Commutator-IO/learn-lean/blob/main/exams/2025/Brevet2025Metropole.lean\#L66}{Brevet2025Metropole.lean\#L66}

\begin{reponse}
Question 5 : en ajoutant une boule \texttt{50} dans chaque urne, l'égalité est rompue \textemdash{}
\texttt{3/7} d'un côté, \texttt{4/10} de l'autre, et \texttt{3/7 > 2/5}.

Ajouter un même cas favorable à deux urnes de tailles différentes ne conserve pas
l'égalité des proportions : c'est l'urne la plus petite qui en profite le plus.
\end{reponse}
\begin{proof}
La boule $50$ est supérieure à $20$ : elle ajoute un cas favorable et un cas possible à
chaque urne. On compte alors trois boules favorables sur sept dans $A$, et quatre sur dix
dans $B$.

Les produits en croix ne sont plus égaux,
\[ 3 \times 10 = 30 \ne 28 = 7 \times 4 , \]
donc $\dfrac{3}{7} \ne \dfrac{4}{10}$ : l'égalité de la question précédente est détruite.
Elle l'est en faveur de la plus petite urne, puisque $\dfrac{3}{7} > \dfrac{4}{10}$.
\end{proof}
\source{https://github.com/Commutator-IO/learn-lean/blob/main/exams/2025/Brevet2025Metropole.lean\#L76}{Brevet2025Metropole.lean\#L76}

\subsection{Exercice 2 \textemdash{} l'aquathlon}

\begin{reponse}
Partie A, question 1 : \texttt{AD = AE − DE = 250 − 50 = 200 m}.

Les points \texttt{A}, \texttt{D}, \texttt{E} sont alignés dans cet ordre : la longueur cherchée est une
différence, ce que la figure donne et que l'énoncé écrit ne dit pas autrement.
\end{reponse}
\begin{proof}
En remplaçant $AE$ par $250$ et $DE$ par $50$ :
\[ AE - DE = 250 - 50 = 200 . \]
\end{proof}
\source{https://github.com/Commutator-IO/learn-lean/blob/main/exams/2025/Brevet2025Metropole.lean\#L87}{Brevet2025Metropole.lean\#L87}

\begin{reponse}
Partie A, question 2 : \texttt{CD = 520 m}.

Le triangle \texttt{ADC} est rectangle en \texttt{A}, donc \texttt{CD\(^{2}\) = AC\(^{2}\) + AD\(^{2}\) = 480\(^{2}\) + 200\(^{2}\)}, soit
\texttt{270 400}, dont la racine est \texttt{520}. Une longueur est positive : c'est ce qui rend la
solution unique.
\end{reponse}

\begin{figure}[h]
\centering
% Construction : A = (0, 0), D = (2, 0), C = (0, 4,8) — une unité vaut 100 m.
% L'angle droit est en A ; l'hypoténuse est [CD], de longueur 520 m.
\begin{tikzpicture}[scale=0.8]
  \coordinate (A) at (0, 0);
  \coordinate (D) at (2, 0);
  \coordinate (C) at (0, 4.8);
  \draw (A) -- (D) -- (C) -- cycle;
  \draw (0.25, 0) -- (0.25, 0.25) -- (0, 0.25);
  \node[below left]  at (A) {$A$};
  \node[below right] at (D) {$D$};
  \node[above left]  at (C) {$C$};
  \node[below] at (1, 0)   {$200$ m};
  \node[left]  at (0, 2.4) {$480$ m};
  \node[above right] at (1, 2.4) {$520$ m};
\end{tikzpicture}

\caption{Le triangle \(ADC\), rectangle en \(A\) : l'hypoténuse \(CD\) se calcule par le théorème de Pythagore.}
\end{figure}
\begin{proof}
Avec $AC = 480$ et $AD = 200$, l'hypothèse devient
\[ CD^2 = 480^2 + 200^2 = 230\,400 + 40\,000 = 270\,400 = 520^2 . \]
Donc $CD^2 - 520^2 = 0$, c'est-à-dire
\[ (CD - 520)(CD + 520) = 0 . \]
Un produit est nul quand l'un de ses facteurs l'est. Comme $CD \ge 0$, le second facteur
vaut au moins $520$, donc il n'est pas nul : c'est le premier qui l'est, et $CD = 520$.

C'est l'hypothèse de positivité qui rend la solution unique \textemdash{} $-520$ vérifie
aussi l'équation du carré, mais n'est pas une longueur.
\end{proof}
\source{https://github.com/Commutator-IO/learn-lean/blob/main/exams/2025/Brevet2025Metropole.lean\#L95}{Brevet2025Metropole.lean\#L95}

\begin{reponse}
Partie A, question 3a : les droites \texttt{(CD)} et \texttt{(BE)} sont parallèles.

Les rapports \texttt{AD/AE} et \texttt{AC/AB} valent tous deux \texttt{0,8} \textemdash{} avec \texttt{AB = AC + CB = 600} \textemdash{}
et les points sont alignés dans le même ordre : la réciproque du théorème de Thalès
conclut. Ce qui se vérifie ici est l'égalité des rapports ; la réciproque elle-même
est un énoncé du chapitre \og{} Géométrie plane \fg{}, et la configuration reste dans la
figure.
\end{reponse}

\begin{figure}[h]
\centering
% Construction : A = (0, 0), C = (4,8 ; 0), B = (6, 0), D = (0, 2), E = (0 ; 2,5).
% Une unité vaut 100 m. A, C, B sont alignés, ainsi que A, D, E ; (CD) et (BE)
% ont la même pente, \(-5/12\).
\begin{tikzpicture}[scale=0.85]
  \coordinate (A) at (0, 0);
  \coordinate (C) at (4.8, 0);
  \coordinate (B) at (6, 0);
  \coordinate (D) at (0, 2);
  \coordinate (E) at (0, 2.5);
  \draw (A) -- (B);
  \draw (A) -- (E);
  \draw (C) -- (D);
  \draw (B) -- (E);
  \draw (0.22, 0) -- (0.22, 0.22) -- (0, 0.22);
  \foreach \p in {A, C, B, D, E} \fill (\p) circle (1.4pt);
  \node[below left]  at (A) {$A$};
  \node[below] at (C) {$C$};
  \node[below] at (B) {$B$};
  \node[left]  at (D) {$D$};
  \node[left]  at (E) {$E$};
  \node[below] at (2.4, -0.28) {$480$ m};
  \node[below] at (5.4, -0.62) {$120$ m};
  \node[left]  at (-0.3, 1) {$200$ m};
\end{tikzpicture}

\caption{La configuration de Thalès du parcours : les rapports \(AD/AE\) et \(AC/AB\) valent tous deux \(0{,}8\).}
\end{figure}
\begin{proof}
Le côté $AB$ vaut $AC + CB = 480 + 120 = 600$. Les deux rapports à comparer sont
$\dfrac{AD}{AE} = \dfrac{200}{250}$ et $\dfrac{AC}{AB} = \dfrac{480}{600}$, et les produits
en croix sont égaux :
\[ 200 \times 600 = 120\,000 = 250 \times 480 . \]
Les deux rapports valent donc tous deux $0{,}8$.

Les points étant alignés dans le même ordre sur les deux demi-droites, la réciproque du
théorème de Thalès permet de conclure au parallélisme de $(CD)$ et $(BE)$. Cette
réciproque est un énoncé du chapitre \og{} Géométrie plane \fg{} ; ce qui est vérifié ici
est l'égalité des rapports, qui en est l'hypothèse.
\end{proof}
\source{https://github.com/Commutator-IO/learn-lean/blob/main/exams/2025/Brevet2025Metropole.lean\#L107}{Brevet2025Metropole.lean\#L107}

\begin{reponse}
Partie B, question 4 : le temps médian des neuf nageurs est \texttt{6 min}, soit
\texttt{360 s} \textemdash{} quatre élèves font mieux, quatre font moins bien.

Les temps sont convertis en secondes, seule façon de les comparer sans ambiguïté :
\texttt{5 min 49 s} et \texttt{5 min 50 s} ne se distinguent pas sur les minutes.
\end{reponse}
\begin{proof}
Les neuf temps, convertis en secondes, sont
\[ 330 ;\ 345 ;\ 349 ;\ 350 ;\ 360 ;\ 371 ;\ 372 ;\ 380 ;\ 400 . \]
On compte ceux qui sont strictement inférieurs à $360$ : $330$, $345$, $349$ et $350$,
soit quatre. Puis ceux qui lui sont strictement supérieurs : $371$, $372$, $380$ et $400$,
soit quatre encore.

L'effectif est impair et se partage en quatre valeurs plus petites, quatre valeurs plus
grandes, et une valeur au milieu : $360$ secondes, c'est-à-dire $6$ minutes, est bien la
médiane.
\end{proof}
\source{https://github.com/Commutator-IO/learn-lean/blob/main/exams/2025/Brevet2025Metropole.lean\#L114}{Brevet2025Metropole.lean\#L114}

\begin{reponse}
Partie B, question 5 : le poisson rouge nage plus vite que le meilleur élève.

Le plus rapide couvre \texttt{200 m} en \texttt{5 min 30 s}, soit \texttt{330 s} : sa vitesse vaut
\texttt{200/330 \(\approx\) 0,61 m/s}, quand \texttt{5 km/h} font \texttt{5 000/3 600 \(\approx\) 1,39 m/s}. Comparer deux
vitesses demande d'abord de les exprimer dans la même unité, et c'est là toute la
question.
\end{reponse}
\begin{proof}
Le meilleur nageur parcourt $200$ mètres en $5$ minutes $30$ secondes, c'est-à-dire
$330$ secondes : sa vitesse est $\dfrac{200}{330}$ mètre par seconde. Le poisson nage à
$5$ kilomètres par heure, soit $\dfrac{5\,000}{3\,600}$ mètre par seconde.

Les deux vitesses sont maintenant dans la même unité, et les produits en croix les
comparent :
\[ 200 \times 3\,600 = 720\,000 < 1\,650\,000 = 330 \times 5\,000 , \]
donc $\dfrac{200}{330} < \dfrac{5\,000}{3\,600}$ : le poisson est le plus rapide.
\end{proof}
\source{https://github.com/Commutator-IO/learn-lean/blob/main/exams/2025/Brevet2025Metropole.lean\#L127}{Brevet2025Metropole.lean\#L127}

\subsection{Exercice 3 \textemdash{} questionnaire à choix multiples}

\begin{reponse}
Question 1 : si \texttt{3} melons coûtent \texttt{8,40 €}, \texttt{5} melons coûtent \texttt{14 €}.

Le prix est proportionnel à la quantité : le prix unitaire est \texttt{2,80 €}.
\end{reponse}
\begin{proof}
Le prix d'un melon est $\dfrac{8{,}40}{3} = 2{,}80$ euro, et cinq melons coûtent
\[ 2{,}80 \times 5 = 14 \text{ euros} . \]
\end{proof}
\source{https://github.com/Commutator-IO/learn-lean/blob/main/exams/2025/Brevet2025Metropole.lean\#L134}{Brevet2025Metropole.lean\#L134}

\begin{reponse}
Question 3 : un article à \texttt{350 €} augmenté de \texttt{20 \%} coûte \texttt{420 €}.

Augmenter de \texttt{20 \%} est multiplier par \texttt{1,2}, non ajouter \texttt{20}.
\end{reponse}
\begin{proof}
Augmenter de $20\,\%$, c'est multiplier par $1 + \dfrac{20}{100} = 1{,}2$ :
\[ 350 \times 1{,}2 = 420 . \]
\end{proof}
\source{https://github.com/Commutator-IO/learn-lean/blob/main/exams/2025/Brevet2025Metropole.lean\#L139}{Brevet2025Metropole.lean\#L139}

\begin{reponse}
Question 4 : l'aire du triangle \texttt{ABC} de côtés \texttt{6}, \texttt{4,5} et \texttt{7,5} vaut
\texttt{13,5 cm\(^{2}\)}.

Le triangle est rectangle en \texttt{B} par la réciproque de Pythagore \textemdash{}
\texttt{6\(^{2}\) + 4,5\(^{2}\) = 56,25 = 7,5\(^{2}\)} \textemdash{}, et son aire est le demi-produit des côtés de l'angle
droit. La réponse \texttt{27} est celle qu'on obtient en oubliant le demi ; la réponse
\texttt{18} celle qu'on obtient en prenant l'hypoténuse pour un côté de l'angle droit.
\end{reponse}
\begin{proof}
D'une part
\[ 6^2 + 4{,}5^2 = 36 + 20{,}25 = 56{,}25 = 7{,}5^2 , \]
donc, par la réciproque du théorème de Pythagore, le triangle est rectangle et son
hypoténuse est le côté de longueur $7{,}5$ \textemdash{} l'angle droit est en $B$.

D'autre part l'aire d'un triangle rectangle est le demi-produit des côtés de l'angle
droit :
\[ \frac{6 \times 4{,}5}{2} = \frac{27}{2} = 13{,}5 . \]
\end{proof}
\source{https://github.com/Commutator-IO/learn-lean/blob/main/exams/2025/Brevet2025Metropole.lean\#L148}{Brevet2025Metropole.lean\#L148}

\begin{reponse}
Question 5 : la forme développée et réduite de \texttt{(2x + 3)(x − 4)} est
\texttt{2x\(^{2}\) − 5x − 12}, pour toute valeur de \texttt{x}.

C'est une identité, non une équation : elle vaut pour tout \texttt{x}, et c'est ce que
\og{} développer \fg{} signifie.
\end{reponse}
\begin{proof}
On développe, en multipliant chaque terme de la première parenthèse par chaque terme de la
seconde :
\[ (2x + 3)(x - 4) = 2x^2 - 8x + 3x - 12 = 2x^2 - 5x - 12 . \]
Le calcul ne suppose rien sur $x$ : l'égalité vaut pour toute valeur, c'est une identité.
\end{proof}
\source{https://github.com/Commutator-IO/learn-lean/blob/main/exams/2025/Brevet2025Metropole.lean\#L156}{Brevet2025Metropole.lean\#L156}

\begin{reponse}
Question 6 : le volume d'une pyramide à base rectangulaire \texttt{7 cm \(\times\) 4 cm} et de
hauteur \texttt{12 cm} vaut \texttt{112 cm\(^{3}\)}.

Le tiers est ce qui distingue la pyramide du prisme de même base et même hauteur :
la réponse \texttt{336} est ce prisme.
\end{reponse}
\begin{proof}
Le volume d'une pyramide est le tiers du produit de l'aire de sa base par sa hauteur :
\[ \frac{7 \times 4 \times 12}{3} = \frac{336}{3} = 112 . \]
\end{proof}
\source{https://github.com/Commutator-IO/learn-lean/blob/main/exams/2025/Brevet2025Metropole.lean\#L164}{Brevet2025Metropole.lean\#L164}

\subsection{Exercice 4 \textemdash{} programmes de calcul}

\begin{reponse}
Partie A, question 1 : le programme de Zoé renvoie \texttt{20} quand on part de \texttt{10}.
\end{reponse}
\begin{proof}
On suit les instructions dans l'ordre :
\[ 2 \times (10 - 4) + 8 = 2 \times 6 + 8 = 12 + 8 = 20 . \]
\end{proof}
\source{https://github.com/Commutator-IO/learn-lean/blob/main/exams/2025/Brevet2025Metropole.lean\#L169}{Brevet2025Metropole.lean\#L169}

\begin{reponse}
Partie A, question 2 : il renvoie \texttt{−14} quand on part de \texttt{−7}.

Le signe du produit est celui qu'il faut surveiller : \texttt{−7 − 4 = −11}, et non \texttt{−3}.
\end{reponse}
\begin{proof}
De même, en surveillant le signe :
\[ 2 \times (-7 - 4) + 8 = 2 \times (-11) + 8 = -22 + 8 = -14 . \]
La soustraction $-7 - 4$ vaut $-11$, et non $-3$.
\end{proof}
\source{https://github.com/Commutator-IO/learn-lean/blob/main/exams/2025/Brevet2025Metropole.lean\#L174}{Brevet2025Metropole.lean\#L174}

\begin{reponse}
Partie A, question 3 : Zoé a raison \textemdash{} son programme renvoie toujours le double
du nombre choisi.

Deux exemples ne prouvent rien ; le calcul littéral, lui, conclut pour tous les
nombres à la fois. C'est le passage du numérique au littéral que la question évalue,
et il est ici l'argument entier.
\end{reponse}
\begin{proof}
On développe et on réduit, en appelant $x$ le nombre choisi :
\[ 2(x - 4) + 8 = 2x - 8 + 8 = 2x . \]
Le calcul vaut pour toute valeur de $x$ : c'est ce qui distingue une démonstration de deux
vérifications sur des exemples.
\end{proof}
\source{https://github.com/Commutator-IO/learn-lean/blob/main/exams/2025/Brevet2025Metropole.lean\#L182}{Brevet2025Metropole.lean\#L182}

\begin{reponse}
Partie B, question 4 : le programme de Fred renvoie \texttt{20x + 50}.

Les trois instructions se composent : \texttt{x \(\mapsto\) 4x \(\mapsto\) 4x + 10 \(\mapsto\) 5(4x + 10)}. L'ordre
compte, et c'est le \texttt{+ 10} avant le \texttt{\(\times\) 5} qui donne le \texttt{+ 50}.
\end{reponse}
\begin{proof}
Les trois instructions se composent : le nombre $x$ devient $4x$, puis $4x + 10$, puis
\[ (4x + 10) \times 5 = 20x + 50 . \]
\end{proof}
\source{https://github.com/Commutator-IO/learn-lean/blob/main/exams/2025/Brevet2025Metropole.lean\#L188}{Brevet2025Metropole.lean\#L188}

\begin{reponse}
Partie B, question 5 : pour obtenir \texttt{75}, il faut choisir \texttt{1,25}.

L'équation \texttt{20x + 50 = 75} a une solution unique, et elle n'est pas entière : le
programme n'est pas magique, mais il est inversible.
\end{reponse}
\begin{proof}
De $20x + 50 = 75$ on tire $20x = 25$, puis
\[ x = \frac{25}{20} = 1{,}25 . \]
\end{proof}
\source{https://github.com/Commutator-IO/learn-lean/blob/main/exams/2025/Brevet2025Metropole.lean\#L194}{Brevet2025Metropole.lean\#L194}

\subsection{Exercice 5 \textemdash{} acheter ou louer}

\begin{reponse}
Partie A, question 1 : avec l'option Achat, la dépense au bout d'un an vaut
\texttt{23 300 €} \textemdash{} le prix de la voiture, plus douze mois d'assurance.
\end{reponse}
\begin{proof}
Une année compte douze mois d'assurance, qui s'ajoutent au prix d'achat :
\[ 22\,400 + 75 \times 12 = 22\,400 + 900 = 23\,300 . \]
\end{proof}
\source{https://github.com/Commutator-IO/learn-lean/blob/main/exams/2025/Brevet2025Metropole.lean\#L201}{Brevet2025Metropole.lean\#L201}

\begin{reponse}
Partie A, question 2 : au bout de \texttt{36} mois, l'option Location fait économiser
\texttt{9 800 €}.

L'écart se creuse d'abord vite \textemdash{} le prix d'achat est payé tout de suite \textemdash{} puis se
resserre, l'assurance étant six fois moins chère que la location.
\end{reponse}
\begin{proof}
Au bout de $36$ mois, l'option Achat coûte $22\,400 + 75 \times 36 = 25\,100$ euros, et
l'option Location $425 \times 36 = 15\,300$ euros. L'écart est
\[ 25\,100 - 15\,300 = 9\,800 . \]
\end{proof}
\source{https://github.com/Commutator-IO/learn-lean/blob/main/exams/2025/Brevet2025Metropole.lean\#L208}{Brevet2025Metropole.lean\#L208}

\begin{reponse}
Partie B, question 5 : l'option Achat devient la plus avantageuse à partir du
soixante-cinquième mois.

L'inéquation \texttt{22 400 + 75x < 425x} se réduit à \texttt{22 400 < 350x}, c'est-à-dire
\texttt{x > 64} : les deux dépenses sont égales au bout de \texttt{64} mois, à \texttt{27 200 €}, et
l'Achat ne devient strictement meilleur qu'ensuite.

Le sujet demande une lecture graphique, et la réponse attendue est \og{} à partir de
\texttt{64} mois \fg{}. Ce qui est démontré ici est l'inéquation exacte, ce que le graphique ne
peut que suggérer \textemdash{} la donnée étant algébrique, contrairement aux lectures de courbe
d'autres sujets.
\end{reponse}

\begin{figure}[h]
\centering
% Construction : abscisse = mois/16, ordonnée = euros/9 000.
% Achat : y = 22 400 + 75x, de (0 ; 2,49) à (96 ; 3,29).
% Location : y = 425x, de (0 ; 0) à (96 ; 4,53).
% Les deux droites se croisent en x = 64 mois, y = 27 200 euros.
\begin{tikzpicture}[xscale=0.85, yscale=0.85]
  \draw[->] (0, 0) -- (6.6, 0) node[below left] {\small mois};
  \draw[->] (0, 0) -- (0, 5.2) node[above left] {\small euros};
  \draw (0, 2.489) -- (6, 3.289);
  \draw[dashed] (0, 0) -- (6, 4.533);
  \draw[dotted] (4, 0) -- (4, 3.022) -- (0, 3.022);
  \fill (4, 3.022) circle (1.6pt);
  \node[below] at (4, -0.05) {\small $64$};
  \node[left]  at (-0.05, 3.022) {\small $27\,200$};
  \node[right] at (6, 3.289) {\small Achat};
  \node[right] at (6, 4.533) {\small Location};
\end{tikzpicture}

\caption{Les deux dépenses en fonction du nombre de mois : elles s'égalisent à \(64\) mois, et l'Achat devient meilleur ensuite.}
\end{figure}
\begin{proof}
Les deux membres de l'inéquation $22\,400 + 75x < 425x$ se réorganisent en retranchant
$75x$ :
\[ 22\,400 < 350x , \]
puis en divisant par $350$, qui est positif et ne change donc pas le sens :
\[ \frac{22\,400}{350} = 64 < x . \]
Chacune de ces deux transformations se défait comme elle s'est faite, ce qui donne les
deux sens de l'équivalence.

Le seuil $64$ n'est pas atteint mais dépassé : au bout de $64$ mois les deux dépenses sont
égales, à $27\,200$ euros, et l'Achat ne devient strictement plus avantageux qu'ensuite.
\end{proof}
\source{https://github.com/Commutator-IO/learn-lean/blob/main/exams/2025/Brevet2025Metropole.lean\#L221}{Brevet2025Metropole.lean\#L221}

\vfill
\noindent\rule{0.35\linewidth}{0.4pt}\par
\noindent{\footnotesize Leonardo de Moura et Sebastian Ullrich,
\emph{The Lean~4 Theorem Prover and Programming Language},
\emph{Automated Deduction — CADE~28}, Springer, 2021, p.~625--635,
\href{https://doi.org/10.1007/978-3-030-79876-5_37}{doi:10.1007/978-3-030-79876-5\_37}.\par}

\end{document}
